\quellseite{10}

\chapter{Widmung an die sächsischen Fürsten}

Denen Durchläuchtigsten, Hochgebornen, Fürsten und Herren, Herrn Friedrichen, Regierenden Landes-Fürsten und Herrn, Herrn Albrechten, Herrn Bernharden, Herrn Heinrichen, Herrn Christianen, Herrn Ernsten, Herrn Johann Ernsten, Sämmtlichen Herren Gebrüdern, Herzogen zu Sachsen, Jülich, Cleve, und Berg, Landgrafen in Thüringen, Markgrafen zu Meißen, Gefürsteten Grafen zu Henneberg, Grafen zu der Mark und Ravensburg, Herren zu Ravenstein Meinen allerseits Gnädigsten Fürsten und Herren.

Durchläuchtigste, Hochgeborne Fürsten, Gnädigste Herren.

Es schreibet der Apostel S. Paulus an die Colosser am 3.: Lasset das Wort Christi reichlich unter euch wohnen in aller Weisheit, lehret und ermahnet euch selbst mit Psalmen und Lobgesängen und geistlichen lieblichen Liedern. Durch das Wort Christi verstehen Lombardus, der Cardinal Cajetanus, die Jesuiten, Cornelius a Lapide, und Eccius, die Lehre, und das Evangelium Christi. Daher auch die Centuriatores Magdeburgenses aus diesen Worten S. Pauli beweisen: die Historia der Evangelisten wäre neben den Büchern Altes Testaments zu der Apostel Zeiten öffentlich, wenn die Gläubigen zusammen kommen, vorgelesen worden. Centur. 1. lib. 2. cap. 6, fol. 494. Der Kirchen-Lehrer Chrysostomus und Theophylactus verstehen durch das Wort Christi, die heilige Schrift, wie bald soll angeführet werden. Diese Lehre nun von Christo, desselben Gnaden-Predigt, oder die heilige Schrift, soll unter denen zu Colossen wohnen. Das ist, diese Lehre soll nicht wie ein Gast mit Freuden angenommen, und hernach bald wieder weggelassen, oder, wie es bei dem ein wenig zu lange bleibenden Gaste meist zu geschehen pfleget, gar verachtet werden, sondern sie soll wohnen, das ist, die Lehre von Christo, die heilige Schrift, soll mit wahrem Glauben ergriffen, und beständig im Herzen bewahret werden, damit wir eine solche Freundschaft mit derselben machen, und sie uns so gemein und bekannt werde, als diejenigen, die bei uns in einem Hause wohnen, uns pflegen gar gemein und kundig zu werden. Und dieses soll geschehen reichlich, das ist, überflüssig und häufig.

Daß also die Art und Weise mit der Lehre Christi, oder der heiligen Schrift umzugehen, hiermit wird angedeutet. Der große majestätische Herr, der Herr der Heerschaaren, befiehlet dem Helden Josua ernstlich: Laß das Buch dieses Gesetzes nicht von deinem Munde kommen, sondern betrachte es Tag und Nacht, auf daß du haltest und thust aller Dinge nach dem, das drinnen geschrieben stehet; im Buch Josua am ersten Kap. v. 8., und vorher sagte er zu allen Israeliten: Diese Worte, die ich dir heute gebiete, sollst du zu Herzen nehmen, und sollst sie deinen Kindern schärfen, und davon reden, wenn du in deinem Hause sitzest, oder auf dem Wege gehest, wenn du dich niederlegest, oder aufstehest, 5 Buch Mosis am 6. Kap. v. 6. König David preiset denjenigen selig, welcher zum Gesetz, zur Lehre des Herrn Lust hat, und redet von seinem Gesetz Tag und Nacht, Psalm 1, v. 2. Unser Heiland Christus Jesus saget zu den Juden: Suchet in der Schrift,

\quellseite{11}

denn ihr meinet, ihr habt das ewige Leben drinnen, und sie ist, die von mir zeuget, Joh. 5, v. 39. Chrysostomus erkläret diese Worte reichlich also: nicht oben hin, sondern mit großem Fleiße.

Was nun Gott selber nennet: Das Gesetz nicht lassen vom Munde kommen; es Tag und Nacht betrachten; zu Herzen nehmen; es den Kindern schärfen; was König David heißt: zum Gesetz des Herrn Lust haben, davon reden Tag und Nacht. Was unser Heiland selber ausspricht mit: Suchen in der Schrift: das redet S. Paulus aus mit: Das Wort Christi reichlich unter sich wohnen lassen. Es wird aber nicht allein mit dem Worte reichlich gesehen auf die Art und Weise, wie man die heilige Schrift solle treiben; sondern auch auf den erbaulichen Nutzen, daß man sie also solle handeln, betrachten und damit umgehen, daß wir reich werden am Glauben und an der Liebe, auch in allen widrigen Begebenissen einen Trost nach dem andern zu unserm, und anderer Leute Besten können alsobald bei der Hand haben. Und dieses reichliche, überflüssige und zu unserer eigenen, auch anderer Leute Erbauung vorgenommenes Treiben der heiligen Schrift soll geschehen unter den Colossern: Unter euch, schreibet S. Paulus, ohne Unterschied des Standes, oder des Geschlechts, es mögen sein Lehrer, oder Zuhörer, Weiber oder Männer. Wie denn der Autor der Erklärungen, welche von etlichen Hieronymo, von andern Primasio, oder sonst jemandem zugereignet werden, diese Worte ausdrücklich von den Laien ausleget, wenn er schreibet: Hic ostenditur, verbum Christi non solum sufficienter, sed etiam abundanter habere debere etiam laicos et invicem semetipsos monere et docere. Das ist: Hier wird gewiesen, daß auch die Laien das Wort Christi nicht allein zur Bedürfniß und Genüge, sondern auch zum Ueberflusse sollen haben, und einer den andern erinnern und unterweisen. Das Chrysostomus unter den Worten: Unter euch, auch die γνωστικος, und βιωτικος, seculares, welche mit zeitlichen Sachen umgehen, in einem weltlichen Stande leben, welche Weib und Kinder haben, begreife, werden wir bald hören.

Dieses reichliche, überflüssige und fleißige Treiben der heiligen Schrift soll ferner geschehen: in aller Weisheit. Der bekannte Schullehrer Thomas, welchen die Römisch-Katholischen insgemein den englischen Lehrer nennen, vermeinet, S. Paulus weise hiermit die Gläubigen zu Colossen auf den rechten Ursprung, woraus alle Weisheit entspringe, und auch auf den Nutzen, welcher von dieser Weisheit zu hoffen, wenn er setzet: Cum dicit: Verbum Christi, monet ad Sapientiam; et primo docet Sapientiae originem, secundo sapientiae usum: ibi, docentes: Das ist: wenn er, Paulus, saget: Das Wort Christi, ermahnet er zur Weisheit; und lehret erstlich der Weisheit Ursprung, zum andern der Weisheit Nutzbarkeit, mit dem Worte: lehret. Wollte man dieses hinzusetzen, daß S. Paulus mit diesen Worten: in aller Weisheit, auch ein Absehen hätte auf die Weise, wie man mit Lesung, Betrachtung und Anführung der heiligen Schrift sollte umgehen, nämlich: behutsam, daß man dasjenige, was vorhergehe, was folge und dergleichen, wohl in Acht nehme; auch zugleich ein Ziel wollte setzen, wie weit man mit Betrachtung der heiligen Schrift gehen sollte, wird es der Aehnlichkeit des Glaubens nicht zuwider laufen. Denn es muß doch mit dem Nachsuchen in der heiligen Schrift Maß gehalten werden, daß wir dieselbe nicht nach unserm Verstande und Gutdünken einrichten, und das öfters verdammliche Quomodo? Wie? damit viele Gottes Allmacht, Weisheit und Wahrheit in Zweifel ziehen, von uns gehöret oder angemerket werden möge.

Endlich wird uns auch vorgestellet die End-Ursache, warum man Gottes Wort, die heilige Schrift, so fleißig treiben, andächtig betrachten und ihm so bekannt und läuffig machen solle. Lehret euch selbst, saget S. Paulus, das ist, unterweiset, unterrichtet euch mit einander. Und auf solche Weise werdet ihr in den Glaubens-Artikeln gegründet und bestätiget werden. Vermahnet euch selbst, schreibet S. Paulus weiter. Das ist: Ist jemand mit einem Fehl übereilet worden, ist er in eine Sünde gefallen, so helfet ihr, die ihr geistlich seid, und nicht aus dem Stande der Gnaden euch durch muthwillige Sünden bringen lassen, ihm wieder zurechte mit sanftmüthigem Geiste, nach S. Pauli Ermahnung in dem 6. Kap. der Epistel an die Galater. Oder wie S. Paulus im 10. Kap. an

\quellseite{12}

die Brüder v. 24. es erkläret: Lasset uns unter einander unser selbst wahrnehmen mit Reizen zur Liebe und guten Werken. Der Kirchenlehrer Chrysostomus, der um das vierhundertste Jahr nach Christi Geburt gelebet, leget vor angeführte Worte S. Pauli, wie sie erkläret worden, auch also aus, welche ich nur in unserer bekannten Muttersprache will anführen. Denn also redet er in der neunten Predigt über das dritte Capitel der Epistel an die Colosser seine Zuhörer zu Constantinopel öffentlich an: Höret zu, so viel euer Laien (βιωτικῶν, victum quaerentes) sein, und dem Weibe und Kindern vorstehet, wie auch euch (Paulus) beschielet, daß ihr vor allem die Schrift leset; und zwar nicht schlecht und oben hin, sondern mit großem Fleiße. Und bald hierauf: Höret alle ihr Laien (βιωτικῶν, victum quaerentes), ich ermahne euch, schaffet euch die Bibel, welche ist eine Arznei der Seelen. Wenn ihr nichts anderes wollet, so schaffet nur das Neue Testament, die Apostelgeschichte, die Evangelia, eure stetigen Lehrmeister. Wenn eine Traurigkeit einfället, so mache dich hierzu, als zu einer Apotheke der Arzneien. Daraus nimm dir einen Trost wider den Unfall. Es mag sich entweder ein Schaden zugetragen haben, oder ein Todesfall, oder ein Verlust der Unsrigen. Ja, mache dich nur nicht allein hierzu, sondern nimm alles an. Behalte es in deinem Verstande und Gedächtniß. Dieses ist die Ursache alles Uebels, daß man die Schrift nicht weiß. Soweit Chrysostomus.

Es wird aber daraus dieses unstreitig geschlossen: Sollen die Laien, diejenigen, welche Weiber und Kinder haben, vor allem die Schrift lesen, sie kaufen, und sich daraus wider alle Unfälle trösten, so muß ja folgen, daß man den Laien die Bibel in ihrer Muttersprache lassen, und sie zur fleißigen Lesung derselben anmahnen müsse und solle. Weil ihnen alles Uebels Ursache ist, die Schrift nicht wissen. Es bemühet sich zwar Joh. Paludanus in seinen also genannten Vindiciis Theologicis adversus verbi Dei Corruptelas part. alter. pag. 276. aus dem Stapletono hierauf zu antworten, wenn er schreibet: Der Apostel lehret deutlich, nicht daß die Schrift solle gelesen werden, sondern daß die Predigt des Worts, die ihnen geschehen, solle wiederholet und erneuert werden, indem sie sich unter einander ermahnen mit Psalmen etc. Denn dieses folget daselbst: Singet dem HErrn in euren Herzen. Dieses aber alles wird deutlich entweder auf die geistlichen Unterredungen unter einander mit den Brüdern, oder auf die besondern, oder Privatandachten gezogen. So weit Paludanus.

Aber vergeblich. Denn S. Paulus setzet zweierlei, welches die Gläubigen zu Colossen thun sollen, wie auch Thomas der Schullehrer angemerkt, und es droben schon angeführet, wenn Paulus spricht: Das Wort Christi, so ermahnet er zur Weisheit, und lehret erstlich den Ursprung der Weisheit, hernach auch den Nutz der Weisheit, in dem Worte „lehre.“ Aber Paludanus menget dieses beides unter einander, als ob es einerlei wäre, das Wort Christi reichlich unter uns wohnen lassen, und dasselbe lehren, oder daraus ermahnen. Die doch also unterschieden sind, als dasjenige, welches etwas wirket, und dasjenige, das da gewirket wird. Denn bei wem das Wort Christi reichlich wohnet, der kann auch andere daraus lehren und ermahnen. Es wird auch Paladano nicht nachgegeben, daß S. Paulus nicht rede von Lesung der Schrift, sondern nur von Wiederholung der Predigt des Worts. Denn wie kann doch das Wort Christi reichlich unter den Colossern wohnen, wenn sie die heilige Schrift nicht haben, und dieselbe nicht lesen? Die Wiederholung der Predigt des Worts gehöret auch mit zum reichlichen Wohnen des Worts Christi unter ihnen; aber es bestehet dieses reichliche Wohnen nicht allein in der Wiederholung des gepredigten Worts, sondern man muß auch die heilige Schrift fleißig und reichlich treiben. Welches nicht geschehen kann, wenn man sie, wie Chrysostomus befiehlt, nicht besitzet und sie nicht unsern stetigen Lehrmeister sein läßt. Welches Lehramt aber die heilige Schrift in einer fremden und unbekannten Sprache nicht kann verrichten, sondern in einer bekannten und auch den Laien läufigen Sprache. Dieses aber kann man Paladano willig nachlassen, daß S. Paulus unter dem reichlichen Wohnen des Worts Christi auch

\quellseite{13}

verstehe die geistlichen Unterredungen unter einander mit den Brüdern, die nach Chrysostomi Meinung sind κοσμικοί, βιωτικοί, seculares, laici, victui quaerendo intenti, Laien, die mit weltlichen Geschäften umgehen, und sich bemühen, wie sie sich und die Ihrigen mit gutem Gewissen hinbringen, oder auch Privatandachten, wenn sie nur aus der heiligen Schrift genommen. Denn solches zu unserm obigen Schlusse genug ist. Der Jesuit Justinianus will der Protestirenden Schluß unvermerket ablehnen, wenn er schreibet über unsern Text: Lasset das Wort Christi etc. Daß nicht jemand meine (welches die Novatores [die Protestirenden damit meinen] unbesonnen erdichten) es müsse die heilige Schrift, als leichte und einem Jeden offenstehend, von eines Jeden Hand angetastet werden, auch nicht einmal die Schuhflicker und Maurer ausgenommen, oder auch nicht, welches gar schön wäre, die Weiberchen: so sehet S. Paulus recht darauf: Lehret und vermahnet euch. Als wollte er sagen: einer unterrichte den andern. Welche aber unser Dolmetscher genennet hat: Vosmet ipsos; hat der Apostel griechisch genennet ἑαυτούς, das ist, euch unter einander, daß nämlich die Verständigeren die Unverständigeren unterweisen sollen. So weit Justinianus.

Aber dieses ist ein wunderlicher Schluß: Einer unterweise den andern. Derowegen soll die heilige Schrift nicht von eines Jeden Hand angetastet werden. Vielmehr gilt der Gegensatz: Einer unterweise den andern. Derowegen soll die heilige Schrift von einem Jeden angegriffen und gelesen werden. Der Schluß ist dieser: welche einander aus der heiligen Schrift sollen unterrichten, denen muß die heilige Schrift zu lesen nachgegeben werden.

Nun aber sollen die κοσμικοί und βιωτικοί, die seculares, welche ihre Nahrung suchen, welche Weiber und Kinder haben, (worunter auch die Gerber, Schuhflicker, Maurer und dergleichen Handwerker ohne einigen Zweifel gehören) einander aus heiliger Schrift unterrichten, derowegen muß diesen benannten Leuten und ihresgleichen nachgelassen werden die heilige Schrift zu lesen.

Der Vorspruch kann nicht geleugnet werden, als von demjenigen, der da wollte befehlen, man sollte mit dem Messer Brod schneiden, und zugleich verbieten, man sollte das Messer, als ein Instrument, dazu nicht brauchen. Man sollte aus einem Brunnen Wasser schöpfen, aber sich hierzu keines Eimers oder einigen anderen Instrumentes bedienen. Denn die heilige Schrift, die da von Gott gegeben, ist nütz zur Lehre, zur Strafe, oder Widerlegung, zur Besserung, zur Züchtigung, oder Unterweisung in der Gerechtigkeit, wie S. Paulus schreibet 2 Tim. 3, v. 16. Wer nun befiehlet, man solle aus der heiligen Schrift unterrichten, aber verbeut die heilige Schrift zu lesen, der befiehlet eine Verrichtung, und untersaget das ordentliche hierzu von GOtt selbst anbefohlene Mittel, die heilige Schrift, die uns doch kann unterweisen zur Seligkeit, durch den Glauben an Christo Jesu, und ist nütze zur Lehre und Unterrichtung. Wer aber also befiehlet und also verbeut, dem ist es mit dem Befehl kein Ernst, sondern spottet nur in dieser höchst wichtigen Sache, welche der Menschen Seligkeit angehet, der elenden Leute noch dazu.

Der Nachspruch ist enthalten in S. Pauli Worten: Lasset das Wort Christi reichlich unter euch wohnen, lehret und vermahnet euch selbst unter einander, und hat denselben der treffliche Lehrer Chrysostomus von denen κοσμικοῖς und βιωτικοῖς verstanden, wie dessen Worte kurz vorher angeführet. Derowegen bleibet auch der Schluß unumgestoßen: Dergleichen Leuten ist nachzulassen, die heilige Schrift in ihrer bekannten Muttersprache zu lesen.

Daß aber Justinianus ferner schreibet: Die Verständigen sollen die Unverständigen unterrichten, daraus folget nicht, daß die heilige Schrift nicht allen zu lesen solle verstattet werden. Denn auch unter den von Justiniano angeführten Schuhflickern und Maurern einer in der heiligen Schrift mehr erfahren als der andere, und also nach Justiniani eigenen Worten, der Verständigere den Unverständigern soll unterrichten. Thomas Stapletonus antwortet auf die Worte Chrysostomi: Ipsum moraliter, non dogmatice haec verba Pauli ad Scripturarum Lectionem traxisse. Aber dieser Unterschied inter moraliter et dogmatice, wie ihn Stapletonus machet, hat keinen Grund in der Theologie. Denn an den-

\quellseite{14}

Lehren hängen gleichsam die Anmahnungen, und folgen aus denselben. Denn wenn die Lehre von GOtt, daß er sei allmächtig, vorgetragen, so kann man daraus ermahnen, ein Jeder wolle in allen Nöthen sein Vertrauen auf diesen allmächtigen GOtt setzen. Wenn die Lehre aus der heiligen Schrift bewiesen, daß GOtt alle Sünde wolle allen bußfertigen Sündern vergeben, so kann man daraus eine Anmahnung ziehen: Alle bußfertigen Sünder wollten bei GOtt Vergebung aller ihrer Sünden suchen, und sich derer gewiß versichern. Und dieses gehet an in allen credendis, in allem, was uns wird vorgeleget aus heiliger Schrift zu glauben. Eben diese Bewandtnis hat es auch mit den faciendis, was GOtt wolle, daß wir thun sollen. Als: wenn der Befehl des wohlthätigen und gütigen GOttes aus dem 58. Kap. Esaia v. 7. kund gemacht: Brich dem Hungrigen dein Brod, und die, so im Elend sind, führe ins Haus. So du einen nacket siehest, so kleide ihn, und entzeuch dich nicht von deinem Fleisch; so folget ohne einigen Zwang die Ermahnung: Man wolle sich der Elenden, Hungrigen und Nackten erbarmen.

Ist derowegen kein Gegensatz zuzulassen, inter moraliter et dogmatice etwas vortragen. Denn weil S. Pauli Lehre ist: Lasset das Wort Christi unter euch reichlich wohnen, so nimmt Chrysostomus aus dieser Lehre die Anmahnung: Ihr Laien, die ihr Weiber und Kinder habt, kaufet euch die Schrift, denn sie ist eurer Seelen Arznei. Nehmet aus derselben wider alle Unfälle Tröstungen, leset die Schrift nicht oben hin, sondern mit großem Fleiße. Denn es ist alles Unglücks eine Ursache, die heilige Schrift nicht wissen. Ja diese Anmahnung ist in obiger Lehre Pauli enthalten und begriffen, daß also fast gar kein Unterschied zu machen in dem Sensu, und in der Meinung, zwischen Pauli Worten: Lasset das Wort Christi unter euch reichlich wohnen; und Chrysostomi Worten: Leset die Schrift nicht oben hin, sondern mit großem Fleiße.

Hat also Chrysostomus aus der Lehre Pauli diese Anmahnung rechtmäßig genommen, und bleibt demnach wahr, was droben geschlossen: Sollen die Laien die Bibel schaffen, und mit großem Fleiße lesen, so muß man ihnen dieselbe in ihrer Muttersprache lassen, und sie anmahnen, solche fleißig zu lesen, und in allen widrigen Begegnissen Tröstungen daraus zu nehmen. Kann man derowegen ohne großen Schrecken, ohne Besorgung göttlichen schweren Zorns, und ohne sonderbare Betrübniß nicht lesen, daß die römischen Bischöfe, welche sich doch vor Christi Statthalter ausgegeben, Pius der Vierte, Sixtus der Fünfte und Clemens der Achte, dieser Namen, sich geweigert, die Sünde zu erlassen denen Laien, welche die Bibel würden lesen. Ja, sie haben den Pfarrern, Beichtvätern und Ordinarien verboten, sie sollten keinen Laien die Absolution sprechen und die Sünde vergeben, wenn sie nicht zuvor die Bibel ausgehändigt. Wie unglücklich aber sind diejenigen Gemeinen, in welchen die Lesung der Bibel den Laien bei Strafe der vorbehaltenen Vergebung der Sünden wird verboten; so glücklich sind die Gemeinen, unter welchen das Wort Christi reichlich wohnet, beides, beim öffentlichen Gottesdienste, wie auch daheim in ihren Häusern, da sie es nicht allein frei und sicher dürfen lesen, sondern auch von ihren Seelsorgern und Lehrmeistern zum freundlichsten dazu angemahnet, und ihnen die großen Nutzbarkeiten, die ihnen daraus zuwachsen, vor Augen gestellet werden.

Dergleichen Glückseligkeit haben auch Durchlauchtigste Fürsten, Gnädigste Herren, E. E. E. E. E. E. Fürstl. Durchl. sämmtliche Fürstenthümer und Länder, daß, GOtt sei ewig dafür gedanket, in Städten und Dörfern das Wort Christi reichlich unter derselben Unterthanen wohnet, die Jugend zur Erlernung des Katechismi, der Psalmen und biblischen Sprüche, wie auch zur Lesung der heiligen Schrift, in den Schulen aller Orten fleißig angeführet, auch die Erwachsenen, daß das Gelernte nicht wieder ausfalle, zur Wiederholung werden angehalten.

Wie nun der auch weiland Durchlauchtigste Fürst und Herr, Herr Ernst, Herzog zu Sachsen, Jülich, Cleve und Berg etc. Euer Fürstl. Durchl. Hochgeehrtester Herr Vater, unsers mehr Christseeligen Andenkens, mein weiland Gnädigster Fürst und Herr, Seine herzliche Freude daran hatte, wenn Ihre Durchl. bei deren Unterthanen die wahre Gottesfurcht konnte befördern; also bezeugeten Hochgedachte Fürstl. Durchl. auch Ihr Gnädigstes Gefallen daran, da Anno 1673 im Juni zu Gotha auf dem Friedenstein, dem Fürstl. Residenz-Hause, unterthänigst ich mein Vorhaben mit Herausgebung einer Bibel eröffnete, indem Sie Ihre Gnädigste Einwilligung

\quellseite{15}

dazu in allen Gnaden gaben, und solche auch durch Dero daselbst wohlbestelltes Consistorium schriftlich auszufertigen Gnädigst befohlen. Welches auch geschehen.

Solche Bibel nun, welche durch hiesiges Waisenhauses Kosten zur Erziehung der Armen darin sich befindlichen Waisen aufgelegt worden, habe E. E. E. E. E. E. Fürstl. Durchlaucht unterthänigst überreichen, und zugleich meinen unterthänigsten Glückwunsch wegen ererbter Fürstenthümer und Landen Altenburgischen Theils, wie auch zur angetretenen Fürstlichen Regierung des Höchsten milden Segen in aller Unterthänigkeit öffentlich wiederholen wollen, meinen GOtt unaufhörlich anflehend, daß Er das Wort Christi reichlich in deren wohllöblichen Fürstenthümern und Landen ferner, und bis an den jüngsten Tag wolle lassen wohnen, und schnell laufen, Euer Fürstl. Durchl. bei glücklicher Regierung und allem Fürstlichen hohen Wohlwesen erhalten, derselben Freunde Freund, und Ihrer Feinde Feind sein und bleiben, auch niemals aus diesem Durchlauchtigsten Hause ermangeln lassen an einem Manne nach seinem Herzen, welcher vor Ihm aus und eingehe!

Welches alles daß es geschehen möge, wünschet von Herzen

\begin{unterschrift}

E. E. E. E. E. E. Fürstl. Durchl.\\Altenburg, den 30. März 1676.\\Unterthänigster Diener und Vorbitter bei GOtt\\Joh. Christfried Sagittarius, Dr.,\\Superint. Gener., Ober-Hofprediger und des Consistorii Assessor.

\end{unterschrift}

\quellseite{16}

\chapter{An den Gottseeligen Leser}

Bei Auflegung dieser Bibel, die mit Herrn \unsicher{Miti} Dieterichs Summarien, wie auch Herrn Francisci Bierlings Vorreden und Beschlüssen (dergleichen sonder Ruhm, noch nicht herauskommen, und also keiner, er sei auch wo oder wer er wolle, mit Recht sich kann beklagen, daß ihm eine einzige Bibel nachgedruckt worden, auch die Begierde, Gottes Ehre zu befördern, sowohl der lieben Stadt Altenburg, als andern evangelischen Städten, nicht unbillig wird zu gönnen sein), habe ich\footnote{1. Vielerley editiones in allerhand Formen mit einander conferiret, aber bei vorgefallenen Zweifeln mich nach derjenigen, welche zu Wittenberg im Jahr Christi 1556 in Folio durch Hans Lufften gedruckt, meistentheils gerichtet.\par 2. Habe ich bei Herrn Dieterichs Summarien dieselbe edition, welche im Jahr Christi 1585 zu Nürnberg in Folio allein, und ohne die biblischen Capitel, in den Druck gegeben, zum Grunde gelegt, und die andern hernach mit diesen Summarien herausgekommenen Bibeln gegen obige gehalten, und mit Verwunderung angemercket, daß in vielen Stellen des wohlverdienten Mannes Worte ausgelassen, und in einem andern, auch öfters widrigen Sinne, befunden worden, daß man mit allem Fug des theuren Lutheri Worte, welche in der Warnung vor der Bibel zu finden, auch allhier kann gebrauchen. Es ist mir oft widerfahren, daß ich der Nachdrucker Druck gelesen, also verfälschet gefunden, daß ich meine eigene Arbeit an vielen Orten nicht erkennet. Wer es nicht will glauben, der lasse die editiones gegen einander halten, da wird es mehr als ihm lieb ist, also befinden.\par 3. Sind die Vorreden und Beschlüsse Herrn Bierlings, welche Anno 1596 zu Breslau, in meiner liebsten Geburtsstadt, der ich alles Gute von Herzen wünsche, in Folio absonderlich herausgegeben, von mir durchgelesen, die Orthographie in Acht genommen, auch die Allegata aus der Bibel aufgeschlagen, und mit Versen an den meisten Orten versehen worden.}

Die Druckfehler, welche über allen angewandten Fleiß mit eingeschlichen, dem Sensui aber nichts benehmen, wolle der christliche und geneigte Leser entweder selbst corrigiren oder doch ins besten vermerken, auch diese Arbeit zu seiner ferneren Erbauung selig gebrauchen, und also das Wort Christi reichlich unter und bei ihm wohnen lassen. Ich begehre für alle aufgewandte Mühe, und darüber mir zugestoßenen Angelegenheiten nicht einen Pfennig, sondern Gott allein soll die Ehre, der christlichen Kirchen die Erbauung, und den hülflosen Waisen allhier der Nutz sein und bleiben. Hiermit sei ein Jeder der göttlichen Gnade zu zeitlicher und ewiger Wohlfahrt treulich befohlen. Altenburg, den 3. April Anno 1676.

\quellseite{17}

\chapter{Vorrede auf das Alte Testament}

Das Alte Testament halten etliche gering, als das dem jüdischen Volke allein gegeben, und nun fort aus sei, und nur von vergangenen Geschichten schreibe, meinen, sie haben genug am Neuen Testament, und geben für, eitel geistliche Sinn im alten Testament zu suchen. Wie auch Origenes, Hieronymus und viel hoher Leute mehr gehalten haben. Aber Christus spricht Joh. 5, 39.: Forschet in der Schrift, denn dieselbige gibt Zeugniß von mir. Und S. Paulus gebeut Timotheo 1 Tim. 4, 13.: Er solle anhalten mit Lesen der Schrift. Und rühmet Röm. 1, 2.: wie das Evangelium sei von Gott in der Schrift verheißen. Und 1 Cor. 15, 3. sagt er: Christus sei, nach Laut der Schrift, von Davids Geblüte kommen, gestorben und vom Tode auferstanden. So weiset uns auch S. Petrus mehr denn einmal hinter in die Schrift, 2 Petr. 1, 19.; 3, 15.

Damit sie uns je lehren, die Schrift des Alten Testaments nicht zu verachten, sondern mit allem Fleiß zu lesen, weil sie selbst das Neue Testament so mächtiglich gründen, und bewähren durchs Alte Testament, und sich darauf berufen. Wie S. Lucas, Apg. 17, 11. schreibet, daß die zu Thessalonich täglich forscheten die Schrift, ob sichs so hielte, wie Paulus lehrte. So wenig nun des Neuen Testaments Grund und Beweisung zu verachten ist, so theuer ist auch das Alte Testament zu achten. Und was ist das Neue Testament anders, denn eine öffentliche Predigt und Verkündigung von Christo, durch die Sprüche im Alten Testament gesetzet, und durch Christum erfüllet?

Daß aber diejenigen, so es nicht besser wissen, eine Anleitung und Unterricht haben, nützlich darinnen zu lesen, habe ich diese Vorrede nach meinem Vermögen, so viel mir Gott gegeben, gestellet. Bitte und warne treulich einen jeglichen frommen Christen, daß er sich nicht stoße an der einfältigen Rede und Geschichte, so ihm oft begegnen wird, sondern zweifle nicht daran, wie schlecht es immer sich ansehen läßt, es seien eitel Wort, Werk, Gericht und Geschicht der hohen göttlichen Majestät, Macht und Weisheit. Denn dies ist die Schrift, die alle Weisen und Klugen zu Narren macht, und allein den Kleinen und Albernen offen stehet, wie Christus sagt Matth. 11, 25.: Darum laß deinen Dünkel und Fühlen fahren, und halte von dieser Schrift, als von dem allerhöchsten, edelsten Heiligthum, als von der allerreichsten Fundgrube, die nimmer genug ausgegründet werden mag, auf daß du die göttliche Weisheit finden mögest, welche Gott hie so albern und schlecht fürgeleget, daß er allen Hochmuth dämpfe. Hier wirst du die Windeln und die Krippe finden, da Christus innen liegt, dahin auch der Engel die Hirten weiset, Luc. 2, 12. Schlechte und geringe Windeln sind es, aber theuer ist der Schatz, Christus, der drinnen liegt.

So wisse nun, daß dies Buch ein Gesetzbuch ist, das da lehret, was man thun und lassen soll, und daneben anzeiget Exempel und Geschichte, wie solche Gesetze gehalten oder übertreten sind. Gleich wie das Neue Testament ein Evangelium oder Gnaden-Buch ist, und lehret, wo mans nehmen soll, daß das Gesetz erfüllet werde. Aber gleich wie im Neuen Testament, neben der Gnaden-Lehre, auch viel andere Lehren gegeben werden, die da Gesetz und Gebot sind, das Fleisch zu regieren, sintemal in diesem Leben der Geist nicht vollkommen wird, noch eitel Gnade regieren kann. Also sind auch im Alten Testament, neben den Gesetzen, etliche Verheißung und Gnaden-Sprüche, damit die heiligen Väter und Propheten unter dem Gesetz im Glauben Christi, wie wir, erhalten sind. Doch wie des Neuen Testaments eigentlich Haupt-Lehre ist, Gnade und Friede durch Vergebung der Sünden in Christo verkündigen, also ist des Alten Testament eigentlich Haupt-Lehre, Gesetz lehren, und Sünde anzeigen, und Gutes fordern, Solches wisse im Alten Testament zu warten.

\quellseite{18}

Und daß wir zuerst auf Moses Bücher kommen, der lehret in seinem ersten Buch, wie alle Creaturen geschaffen sind, und, daß seines Schreibens meiste Ursache ist, wo die Sünde und der Tod herkommen sey, nämlich, durch Adams Fall, aus des Teufels Bosheit, 1 Mos. 3, 4. Aber bald darauf, ehe denn Moses Gesetz kömmt, lehret er, woher die Hülfe wiederkommen sollt, die Sünde und Tod zu vertreiben, nämlich, nicht durch Gesetz noch eigen Werk, weil noch kein Gesetz war, sondern durch des Weibes Samen, Christum, Adam, 1 Mos. 3, 15. und Abraham, 1 Mos. 12, 3. verheißen. Auf daß also der Glaube von Anfang der Schrift durch und durch gepreiset werde über alle Werke, Gesetz und Verdienst. Also hat das erste Buch Mose fast eitel Exempel des Glaubens und Unglaubens, und was Glaube und Unglaube für Früchte tragen, und ist fast ein evangelisch Buch.

Danach im andern Buch, da die Welt nun voll in der Blindheit versunken war, daß man schier nicht mehr wußte, was Sünde war, oder wo der Tod herkommen sey, bringet Gott Mosen herfür mit dem Gesetz, und nimmt ein besonders Volk an, die Welt an ihnen wieder zu erleuchten, und durchs Gesetz die Sünde wieder zu eröffnen. Und verfaßt also das Volk mit allerlei Gesetzen, und sondert sie von allen andern Völkern, lässet sie eine Hütten bauen, und richtet einen Gottesdienst an, bestellet Fürsten und Amtleute, und versorget also sein Volk, beide mit Gesetzen und Leuten, aufs allerfeinste, wie sie beide leiblich für der Welt, und geistlich für Gott sollen regieret werden.

Im dritten Buch wird insonderheit das Priesterthum verordnet mit seinen Gesetzen und Rechten, darnach die Priester thun, und das Volk lehren sollen. Da stehet man, wie ein priesterlich Amt nur um der Sünde willen wird eingesetzt, daß es dieselbige soll dem Volk kund machen, und für Gott versöhnen. Also, daß alle sein Werk ist mit Sünden und Sündern umgehen. Derhalben auch den Priestern kein zeitlich Gut gegeben, noch leiblich zu regieren befohlen oder zugelassen wird, sondern allein des Volkes zu pflegen in den Sünden, ihnen zugeeignet wird.

Im vierten Buch, da nun die Gesetze gegeben, Priester und Fürsten eingesetzt sind, die Hütten und Gottesdienst angerichtet, und alles bereit ist, was zum Volk Gottes gehört, hebt sich das Werk und Uebung an, und wird versucht, wie solche Ordnung gehen und sich schicken will. Darum schreibet dasselbige Buch von so viel Ungehorsam und Plagen des Volkes. Und werden etliche Gesetze verkläret und gemehret. Denn also findet sichs allezeit, daß Gesetze bald zu geben sind; aber wenn sie sollen angehen, und in den Schwang kommen, da begegnet nicht mehr denn eitel Hinderniß, und will niemand fort, wie das Gesetz fordert, daß dies Buch ein merklich Exempel ist, wie gar es nichts ist, mit Gesetzen die Leute fromm machen, sondern wie S. Paulus sagt: Daß Gesetz nur Sünde und Zorn anrichte, Röm. 4, 15.

Im fünften Buch, da nun das Volk um seinen Ungehorsam gestraft ist, und GOtt sie mit Gnaden ein wenig gelocket hatte, daß sie aus Wohlthat, da er ihnen die zwei Königreiche gab, beweget wurden, seine Gesetze mit Lust und Liebe zu halten, wiederholt Mose das ganze Gesetz mit allen Geschichten, so ihnen begegnet war (ohne was das Priesterthum betrifft), und erkläret also von neuem an alles, was beide zum leiblichen, und geistlichen Regiment eines Volks gehöret. Daß also Mose, wie ein vollkommener Gesetzlehrer, allenthalben seinem Amt genug thäte, und das Gesetz nicht allein gäbe, sondern auch dabei wäre, da mans thun sollt, und wo es fehlet, verkläret, und wieder anrichtet. Aber diese Verklärung im fünften Buch hält eigentlich nichts anders innen, denn den Glauben an GOtt, und Liebe zum Nächsten. Denn dahin gelangen alle Gesetze Gottes. Darum wehret Mose mit seinem Erklären alle dem, daß den Glauben an GOtt verderben mag, bis hinan in das zwanzigste Kapitel, und alle dem, daß die Liebe hindert, bis an des Buches Ende.

Hierbei ist nun zu merken aufs Erste, daß Mose das Volk so genau mit Gesetzen verfaßt, daß er keinen Raum lässet der Vernunft irgend ein Werk zu erwählen oder eigenen Gottesdienst erfinden. Denn er lehret nicht allein GOtt fürchten, trauen und lieben, sondern gibt auch so mancherlei Weise äußerlichen Gottesdienstes, mit Opfern, Geloben, Fasten, Kasteien \&c., daß niemand noth sey.

\quellseite{19}

etwas anders zu erwählen. Item, er lehret auch pflanzen, bauen, freien, streiten, Kinder, Gesinde und Haus regieren, kaufen und verkaufen, borgen und lösen, und alles was äußerlich und innerlich zu thun sei, sogar, daß etliche Satzungen gleich närrisch und vergeblich anzusehen sind.

Lieber, warum thut GOtt das? Einlich darum, er hat sich des Volks unterwunden, daß es sein eigen sein sollte, und er wollte ihr GOtt sein, darum wollt er sie also regieren, daß alle ihr Thun gewiß wäre, daß es für ihn recht wäre. Denn wo jemand etwas thut, da Gottes Wort nicht zuvor auf gegeben ist, das gilt für GOtt nicht, und ist verloren. Denn er verbeut auch im 5. Buch, Kap. 4, 2., Kap. 12, 32., daß sie nichts sollten zuthun zu seinen Gesetzen. Und im Kap. 12, 8. spricht er: Sie sollen nicht thun, was sie recht dünket. Auch der Psalter und alle Propheten drob schreien, daß das Volk gute Werke thäte, die sie selbst erwähleten, und von GOtt nicht geboten wären. Denn er will und kanns nicht leiden, daß die Seinen etwas fürnehmen zu thun, daß er nicht befohlen hat, es sei wie gut es immer sein kann. Denn Gehorsam ist aller Werk Adel und Güte, der an Gottes Wort hanget.

Weil denn nun dies Leben nicht kann ohn äußerlichen Gottesdienst und Weise sein, hat er ihnen fürgeleget solche mancherlei Weise, und mit seinem Gebot verfaßet, auf daß, ob sie ja müßten oder auch wollten GOtt irgend einen äußerlichen Dienst thun, daß sie dieser einen angreifen, und nicht einen eigenen erdächten, damit sie gewiß und sicher wären, daß solch ihr Werk in Gottes Wort und Gehorsam ginge. Also ist ihnen allenthalben gewehret, eigener Vernunft und freiem Willen zu folgen, Gutes zu thun, und wohl zu leben, und doch übrig genug Raum, Stätte, Zeit, Person, Werk und Weise bestimmt, und fürleget, daß sie nicht klagen dürfen noch fremder Gottesdienst Exempel nachfolgen müssen.

Aufs Andere ist zu merken, daß die Gesetze dreierlei Art sind. Etliche, die nur von zeitlichen Gütern sagen, wie bei uns die käiserlichen Gesetze thun. Diese sind von GOtt allgemein um der Bösen willen gesetzt, daß sie nichts ärgers thäten. Darum sind solche Gesetze nur Wehr-Gesetze, mehr denn Lehr-Gesetze. Als da Mose gehet, ein Weib mit einem Scheidebrief von sich zu lassen, 5 Mos. 24. 1. Item, daß ein Mann sein Weib mit einem Eiferopfer treiben, 4 Mos. 5, 14. und andere Weiber mehr nehmen mag, 5 Mos. 25, 5.; solches sind alles weltliche Gesetze. Etliche aber sind, die von äußerlichem Gottesdienst lehren, wie droben gesagt ist.

Aber diese beide gehen nun die Gesetze vom Glauben und von der Liebe, also, daß alle andern Gesetze müssen und sollen ihr Maß haben vom Glauben und von der Liebe, daß sie gehen sollen, wo ihre Werke also geraten, daß sie nicht wider den Glauben und die Liebe gehen, wo sie aber wider den Glauben und die Liebe geraten, sollen sie schlecht ab sein.

Daher lesen wir, daß David den Mörder Joab nicht tödtet, so er doch zweimal den Tod verdienet hatte, 2 Sam. 3, 27., 20, 10., und 2 Sam. 14, 11., gelobet er dem Weibe von Thekoa, ihr Sohn solle nicht sterben, ob er wohl seinen Bruder erwürget hätte. Item, Absalom tödtet er auch nicht, 2 Sam. 14, 21. Item, er selbst David aß von dem heiligen Brod der Priester, 1 Sam. 21, 6. Item, Thamar meinet, der König möchte sie geben Amnon ihrem Stiefbruder zur Ehe, 2 Sam. 13, 13. Aus diesen und dergleichen Geschichten siehet man wohl, daß die Könige, Priester und Obersten haben oft frisch ins Gesetz gegriffen, wo es der Glaube und die Liebe haben gefordert. Daß also der Glaube und die Liebe soll aller Gesetze Meisterin sein, und sie alle in ihrer Macht haben. Denn sintemal alle Gesetze auf den Glauben und Liebe treiben, soll keines nicht mehr gelten, noch ein Gesetz sein, wo es dem Glauben oder der Liebe will zuwider geraten.

Derhalben irren die Juden noch heutiges Tages fast sehr, daß sie so strenge und hart über etlichen Gesetzen Mose halten, und viel ehe Liebe und Friede lassen untergehen, ehe sie mit uns essen und trinken, oder dergleichen thäten, und sehen des Gesetzes Meinung nicht recht an. Denn dieser Verstand ist vonnöthen allen, die unter Gesetzen leben, nicht allein den Juden. Denn also saget auch Christus Matth. 12, 11., daß man den Sabbath brechen möchte, wo ein Ochs in eine Grube gefallen wäre, und ihm heraus helfen, welches doch nur eine zeitliche Noth und Schaden war. Wie vielmehr soll man frisch allerlei Gesetze brechen, wo es Liebes-Noth fordert, so anders dem Glauben und der Liebe nichts zuwider geschieht. Wie Christus sagt, daß David gethan hat, da er die heiligen Brode aß, Matth. 12, 4.

\quellseite{20}

Was ist aber, daß Mose die Gesetze so unordentlich untereinander wirft? Warum setzet er nicht die weltlichen auf einen Haufen, die geistlichen auch auf einen Haufen, und den Glauben und Liebe auf einen? Dazu wiederholet er zuweilen ein Gesetz so oft und treibet einerlei Wort so viel mal, daß gleich verdrossen ist zu lesen und zu hören? Antwort: Mose schreibt, wie sich’s treibet, daß sein Buch ein Bild und Exempel ist des Regiments und Lebens; denn also gehet es zu, wenn es im Schwang gehet, daß jetzt dies Werk, jetzt jenes gethan sein muß. Und kein Mensch sein Leben also fassen mag (so es anders göttlich sein soll), daß er diesen Tag eitel geistliche, den andern eitel weltliche Gesetze übe, sondern Gott regieret also alle Gesetze unter einander, wie die Sterne am Himmel, und die Blumen auf dem Felde stehen, daß der Mensch muß alle Stunde zu jeglichem bereit sein, und thun, welches ihm am ersten für die Hand kömmt: also ist Mose Buch auch untereinander gemenget.

Daß er aber so fast treibet, und oft einerlei wiederholet, da ist auch seines Amts Art angezeiget. Denn wer ein Gesetzvolk regieren soll, der muß immer anhalten, immer treiben, und sich mit dem Volk wie mit Eseln bläuen. Denn kein Gesetzwerk gehet mit Lust und Liebe abe, es ist alles erzwungen und abgenöthiget. Weil nun Mose ein Gesetzlehrer ist, muß er mit seinem Treiben anzeigen, wie Gesetzwerke gezwungene Werke sind, und das Volk müde machen, bis es durch solch Treiben erkenne seine Krankheit und Unlust zu Gottes Gesetz, und nach der Gnade trachte, wie folget.

Als Dritte ist das die rechte Meinung Mose, daß er durchs Gesetz die Sünde offenbare, und alle Vermessenheit menschlichen Vermögens zu Schanden mache. Denn daher nennet ihn S. Paulus, Gal. 2, 19. einen Amtmann der Sünde, und sein Amt ein Amt des Todes, 2 Cor. 3, 6. und Röm. 3, 20. und Kap. 7, 7. spricht er: Durchs Gesetz komme nicht mehr denn Erkenntniß der Sünde. Und Röm. 3, 20.: Durch Gesetzes Werke wird niemand fromm für Gott. Denn Mose kann durchs Gesetz nicht mehr thun, weder anzeigen, was man thun und lassen soll. Aber Kraft und Vermögen solches zu thun und zu lassen, gibt er nicht, und läßt uns also in der Sünde stecken. Wenn wir denn in der Sünde stecken, so dringet der Tod alsbald auf uns, als eine Rache und Strafe über die Sünde. Daher nennet S. Paulus die Sünde des Todes Stachel 1 Cor. 15, 56., daß der Tod durch die Sünde alle sein Recht und Macht an uns hat. Aber wo das Gesetz nicht wäre, so wäre auch keine Sünde, Röm. 4, 15. Darum ist’s alles Mose Amt Schuld, der reget und rüget die Sünde durchs Gesetz, so folget der Tod auf die Sünde mit Gewalt. Daß Mose Amt billig und recht ein Amt der Sünde und des Todes von S. Paulo genennet wird, 2 Cor. 3, 6.; denn er bringet nichts auf uns durch sein Gesetzgeben, denn Sünde und Tod.

Aber doch ist solches Sündenamt und Todesamt gut und fast von nöthen: denn wo Gottes Gesetz nicht ist, da ist alle menschliche Vernunft so blind, daß sie die Sünde nicht mag erkennen, Röm. 3, 20., Kap. 7, 7. 8., denn keine menschliche Vernunft weiß, daß Unglaube und an Gott verzweifeln Sünde sei, ja sie weiß nichts davon, daß man Gott glauben und trauen soll. Gehet also dahin in ihrer Blindheit verstocket, und fühlet solche Sünde nimmermehr. Thut dieweil sonst etwa gute Werke, und führet ein äußerlich ehrbares Leben. Da meinet sie denn, sie stehe wohl, und sei der Sachen genug geschehen. Wie wir sehen an den Heiden und Heuchlern, wenn sie auf ihr bestes leben. Item, so weiß sie auch nicht, daß böse Neigung des Fleisches und Haß wider die Feinde Sünde sei; sondern weil sie stehet und fühlet, daß alle Menschen so geschickt sind, achtet sie solches für natürlich und recht gut Ding, und meinet, es sei genug, wenn man nur äußerlich den Werken wehret. Also gehet sie dahin, und achtet ihre Krankheit für Stärke, ihre Sünde für recht, ihr Böses für gut, und kann nicht weiter.

Siehe, diese Blindheit und verstockte Vermessenheit zu vertreiben, ist Mose Amt noth. Nun kann er sie nicht vertreiben, er muß sie offenbaren und zu erkennen geben. Das thut er durchs Gesetz, da er lehret, man solle Gott fürchten, trauen, glauben und lieben, dazu keine böse Lust noch Haß zu einigem Menschen tragen oder haben. Wenn nun die Natur solches Recht hört, so muß sie erschrecken. Denn sie befindet gewiß weder Trauen noch Glauben, weder Furcht noch Liebe zu Gott, item, weder Liebe noch Reinigkeit gegen den Nächsten, sondern eitel Unglauben, Zweifeln, Verachtung und Haß zu Gott, und eitel bösen Willen und Lust zum

\quellseite{21}

Nächsten. Wenn sie aber solches findet, so ist der Tod alsobald vor Augen, der solchen Sünder fressen, und in die Hölle will verschlingen.

Siehe, daß heißt den Tod durch die Sünde auf uns bringen, und durch die Sünde uns tödten. Das heißt durch das Gesetz die Sünde regen, und für die Augen setzen, und alle unsere Vermessenheit in ein Verzagen und Zittern und Verzweifeln treiben, daß der Mensch nicht mehr kann thun, denn mit dem Propheten schreien: Ich bin von GOtt verworfen, Esaias 49, 14. oder wie man auf deutsch saget: Ich bin des Teufels, ich kann nimmermehr selig werden. Das heißt recht in die Hölle geführet. Das meinet S. Paulus mit kurzen Worten, 1 Cor. 15, 15., der Stachel des Todes ist die Sünde, aber das Gesetz ist der Sünden Kraft. Als sollte er sagen: Daß der Tod sticht und uns erwürget, machet die Sünde, die an uns erfunden wird, des Todes schuldig. Daß aber die Sünde an uns funden wird, und so mächtig uns dem Tod gibt, machet das Gesetz, welches uns die Sünde offenbaret und erkennen lehret, die wir zuvor nicht kannten und sicher waren.

Nun siehe, mit welcher Gewalt Mosis solches sein Amt treibet und ausrichtet. Denn daß er ja die Natur aufs allerhöchste schände, gibt er nicht allein solche Gesetze, die von natürlichen und wahrhaftigen Sünden sagen, als da sind die zehen Gebote, sondern machet auch Sünde, da von Natur sonst keine Sünde ist, und dringet und drücket auf sie mit Haufen Sünden. Denn Unglaube und böse Lust ist von Art Sünde, und des Todes werth. Aber daß man nicht soll gesäuert Brod essen auf Ostern, 2 Mos. 12, 15. und kein unrein Thier essen, 3 Mos. 11, 4. u. folg. sein Zeichen an dem Leib machen, 3 Mos. 19, 28. 21, 5. und alles, was das Levitische Priesterthum mit Sündern schaffet, das ist nicht von Art Sünde und böse, sondern wird allein darum Sünde, daß durchs Gesetz verboten ist, welches Gesetz wohl kann ab sein, aber die zehen Gebote mögen nicht also ab sein. Denn da ist Sünde, obschon die Gebote nicht wären, oder nicht erkennet wären. Gleichwie der Heiden Unglaube Sünde ist, ob sie es wohl nicht wissen noch achten, daß Sünde sei.

Also sehen wir, daß solche und so mancherlei Gesetze Mose nicht allein darum gegeben sind, daß niemand etwas eigenes dürfte erwählen Gutes zu thun, und wohl zu leben, wie droben gesagt ist, sondern vielmehr darum, daß der Sünden nur viel würden, und sich über die Maß häufeten, daß Gewissen zu beschweren, auf daß die verstockte Blindheit sich erkennen müßte, und ihr eigen Unvermögen und Nichtigkeit zum Guten müßte fühlen, und also durchs Gesetz genöthiget und gedrungen würde etwas weiters zu suchen, denn das Gesetz und eigen Vermögen, nämlich, Gottes Gnade in künftigen Christum verheißen. Denn es ist je alles Gesetz Gottes gut und recht, wenn er auch gleich hieße nur Mist tragen, oder Strohhulm aufheben. So muß aber der ja nicht fromm noch gutes Herzens sein, der solch gut Gesetz nicht hält, oder ungern hält. So vermag alle Natur nichts anders, denn ungerne halten, darum mußte sie hie am guten Gesetz Gottes ihre Bosheit erkennen und fühlen, und nach der Hülfe göttlicher Gnaden seufzen und trachten in Christo.

Darum, wo nun Christus kommt, da höret das Gesetz auf, sonderlich das Levitische, welches Sünde machet, da sonst von Art keine Sünde ist, wie gesagt ist. So hören auch die zehn Gebote auf, nicht also, daß man sie nicht halten noch erfüllen sollte, sondern Mosis Amt höret drinnen auf, daß es nicht mehr durch die zehn Gebote die Sünde stark machet, und die Sünde nicht mehr des Todes Stachel ist. Denn durch Christum ist die Sünde vergeben, GOtt versöhnet, und das Herz hat angefangen dem Gesetz hold zu sein, daß es Mosis Amt nicht mehr kann strafen, und zu Sünden machen, als hätte es die Gebote nicht gehalten, und wäre des Todes schuldig, wie es wohl vor der Gnade, und ehe denn Christus da war.

Das lehret S. Paulus, 2 Cor. 3, 14., da er spricht: daß die Klarheit im Angesicht Mose aufhöret, um der Klarheit willen im Angesicht Jesu Christi. Das ist das Amt Mose, das uns zu Sünden und Schanden machet, mit dem Glanz der Erkenntniß unserer Bosheit und Nichtigkeit, thut uns nicht mehr weh, schrecket uns auch nicht mehr mit dem Tode. Denn wir haben nun die Klarheit im Angesicht Christi, das ist, das Amt der Gnaden, dadurch wir Christum erkennen, mit welches Gerechtigkeit, Leben und Stärke wir das Gesetz erfüllen, Tod und Hölle

\quellseite{22}

überwinden. Wie auch die drei Apostel auf dem Berge Thabor Mosen und Eliam sahen, und doch nicht vor ihnen erschraken, um der lieblichen Klarheit willen im Angesichte Christi, Matth. 17, 4., Luc. 9, 33. Aber da Christus nicht gegenwärtig war, konnten die Kinder Israel die Klarheit und Glanz in Mose Angesicht nicht erleiden, darum mußte er eine Decke dafür thun, 2 Mos. 34, 30. 33.

Denn es sind dreierlei Schüler des Gesetzes: die ersten, die das Gesetz hören und verachten, führen ein ruchlos Leben, ohne Furcht, zu diesen kommt das Gesetz nicht. Und sind bedeutet durch die Kalbthiener in der Wüsten, um welcher willen Mose die Tafeln entzwei warf, und das Gesetz nicht zu ihnen brachte, 2 Mos. 32, 6. 19.

Die andern, die es angreifen mit eigener Kraft zu erfüllen, ohne Gnade, die sind bedeutet durch die, so Mose Antlitz nicht sehen konnten, da er zum andern mal die Tafeln brachte, 2 Mos. 34, 30., zu diesen kommt das Gesetz, aber sie leidens nicht, darum machen sie eine Decke darüber, und führen ein heuchelisch Leben mit äußerlichen Werken des Gesetzes, welches doch das Gesetz alles zu Sünden macht, wo die Decke abgethan wird. Denn das Gesetz erweiset, daß unser Vermögen nichts sei ohne Christus Gnade.

Die dritten sind, die Mosen klar ohne Decke sehen. Das sind die, die des Gesetzes Meinung verstehen, wie es unmöglich Ding fordere. Da gehet die Sünde in der Kraft, 1 Cor. 15, 56., da ist der Tod mächtig, da ist des Goliaths Spieß wie ein Weberbaum, und sein Stachel hat sechshundert Seckel Erz, daß alle Kinder Israel für ihm fliehen, ohne der einige David Christus unser Herr, erlöset uns von dem allen. Denn wo nicht Christus Klarheit neben solcher Klarheit Mose käme, könnte niemand solche Glänze des Gesetzes, der Sünde und des Todes Schrecken ertragen: diese fallen ab von allen Werken und Vermessenheit, und lernen am Gesetz nicht mehr, denn allein Sünde erkennen, und nach Christo zu seufzen. Welches auch das eigentliche Amt Mose und des Gesetzes Art ist, Gal. 3, 24.

Also hat Mose auch selbst angezeigt, daß sein Amt und Lehre soll währen bis auf Christum, und alsdenn aufhören, da er spricht, 5 Mos. 18, 15.: Einen Propheten wird dir der Herr dein GOTT erwecken aus deinen Brüdern, wie mich, den sollst du hören, etc. Dies ist der edelste Spruch, und freilich der Herr im ganzen Mose, welchen auch die Apostel hoch geführet und stark gebraucht haben, Joh. 1, 45., Kap. 6, 14.; Apost. 3, 22., Kap. 7, 37., das Evangelium zu bekräftigen, und das Gesetz abzuthun, und alle Propheten gar viel daraus gezogen. Denn weil Gott hier einen andern Propheten verheißet, den sie hören sollen, zwinget sichs, daß er etwas anders lehren würde, denn Mose, und Mose seine Macht ihm übergibt und weichet, daß man jenen hören solle; so kann jener Prophet nicht Gesetz lehren, denn das hat Mose aufs allerhöchste ausgerichtet, und wäre keine Noth um Gesetzes willen einen andern Propheten zu erwecken, darum ists gewiß von der Gnadenlehre und Christo gesagt.

Darum nennet auch S. Paulus Mose Gesetz das Alte Testament, 2 Cor. 3, 14., Christus auch, da er darum ein Testament eingesetzt, Matth. 26, 28.; und ist darum ein Testament, daß Gott darinnen verhieß und beschied dem Volk Israel das Land Canaan, wo sie es halten würden; und gab's ihnen, und ward bestätiget durch Schöps- und Bocks-Tod und Blut, 2 Mos. 24, 8. Aber weil solch Testament nicht auf Gottes Gnaden, sondern auf Menschen-Werken stund, mußte es alt werden und aufhören, und das verheißene Land wieder verloren werden, darum, daß durch Werke das Gesetz nicht kann erfüllet werden. Und mußte ein ander Testament kommen, das nicht alt würde, auch nicht auf unserm Thun, sondern auf Gottes Wort und Werken stände, auf daß es ewiglich währete, Hebr. 13, 20. Darum ists auch durch einer ewigen Person Tod und Blut bestätiget, und ein ewiges Land verheißen und gegeben. Das sei nun von Mose Büchern und Amt geredet.

Was sind aber nun die andern Bücher der Propheten und der Geschichten? Antwort: Nichts anders, denn was Mose ist, denn sie treiben allesamt Moses Amt, und wehren den falschen Propheten, daß sie das Volk nicht auf die Werke führen, sondern in dem rechten Amt Mose und Erkenntniß des Gesetzes bleiben lassen. Und halten fest darauf, daß sie durch des Gesetzes rechten Verstand die Leute in ihrer eigenen Untüchtigkeit behalten, und auf Christum treiben, wie Mose thut. Darum streichen sie auch weiter aus, was Mose von Christo gesagt hat, und zeigen an

\quellseite{23}

beiderlei Exempel, derer die Mose recht haben, und derer die ihn nicht recht haben, und aller beider Strafe und Lohn. Also, daß die Propheten nichts anders sind, denn Handhaber und Zeugen Mose und seines Amts, daß sie durchs Gesetz jedermann zu Christo bringen.

Aufs letzte sollte ich auch wohl die geistliche Deutung anzeigen, so durch das Levitische Gesetz und Priestertum Mose fürgelegt. Aber es ist sein zuviel schreiben, es will Raum und Zeit haben, und mit lebendiger Stimme ausgelegt sein. Denn freilich Mose ein Brunn ist aller Weisheit und Verstandes, daraus gequollen ist alles, was alle Propheten gewußt und gesagt haben, dazu auch das Neue Testament heraus fleußt, und darin gegründet ist, wie wir gehöret haben. Aber doch ein kleines kurzes Grifflein zu geben, denjenigen, so Gnade und Verstand haben, weiter darnach zu trachten, sei das mein Dienst.

Wenn du willst wohl und sicher deuten, so nimm Christum für dich, denn das ist der Mann, dem es alles ganz und gar gilt. So mache nun aus dem Hohenpriester Aaron niemand denn Christum alleine, wie die Epistel an die Hebräer thut, Kap. 7, 15 folg., welche fast allein genugsam ist, alle Figuren Mose zu deuten. Also ists auch gewiß, daß Christus selbst das Opfer ist, ja auch der Altar, der sich selbst mit seinem eigenen Blut geopfert hat. Wie auch dieselbe Epistel meldet, Kap. 7, 27., Kap. 13, 10. Wie nun der Levitische Hohepriester durch solch Opfer nur die gemachten Sünden wegnahm, die von Natur nicht Sünde waren: also hat unser Hohepriester Christus, durch sein selbst Opfer und Blut die rechte Sünde, die von Natur Sünde ist, weggenommen. Und ist einmal durch den Vorhang gegangen zu Gott, daß er uns versühne. Also, daß du alles, was vom Hohenpriester geschrieben ist, auf Christum persönlich, und sonst auf niemand deutest.

Aber des Hohenpriesters Söhne, die mit dem täglichen Opfer umgehen, sollten auf uns Christen deuten, die wir für unsern Vater Christo im Himmel sitzend, hie auf Erden mit dem Leibe wohnen, und nicht hindurch sind bei ihm, ohne mit dem Glauben geistlich. Derselben Amt, wie sie schlachten und opfern, bedeutet nichts anders, denn das Evangelium predigen, durch welches der alte Mensch getödtet, und Gott geopfert durchs Feuer der Liebe, im heiligen Geist verbrannt und verzehret wird, welches gar wohl reucht für Gott, das ist, es macht ein gut, rein, sicher Gewissen für Gott. Diese Deutung trifft S. Paulus, Röm. 12, 1., da er lehret, wie wir unsere Leiber sollen opfern Gott zum lebendigen, heiligen, angenehmen Opfer. Welches wir thun (wie gesagt) durch stetige Uebung des Evangelii, beide mit Predigen und Glauben. Das sei diesmal genug zur kurzen Anleitung, Christum und das Evangelium zu suchen im Alten Testament.

\quellseite{24}

\chapter{Vitus Dieterich, an den christlichen Leser}

Ich bitte alle frommen Christen, die diese Summaria lesen, sie wollen meine Arbeit und Dienst ihnen gefallen lassen, und zum besten deuten. Denn ichs je christlich gemeinet, und ferner nicht gedacht habe, denn wie durch kurzen und einfältigen Unterricht auch die Ungelehrten möchten zur heiligen Schrift recht geleitet werden. Derhalben ich die scharfe und subtile Auslegung mit Fleiß gemieden, und allenthalben mich an dem habe genügen lassen, daß verständlich, einfältig, und der eigentlichen Deutung der Schrift gemäß gewesen. So hat es sich keineswegs schicken wollen, daß ich es hätte weitläufiger gemacht, weil ich anfänglich diese Arbeit für mich genommen habe, daß solche Summaria des Alten Testaments in meiner Kirchen, vor den Kapiteln gelesen würden, und jetzund viel andere Kirchen solchen Brauch auch angenommen haben, und diese Summaria zum kurzen Unterricht lesen. Denn es muß zwischen solcher Lection und der Predigt ein Unterschied sein. In der Predigt soll man genau allem nachgehen, und nichts dahinten lassen. In solcher Lection aber ists genug, daß man nur auf die vornehmsten Punkte sehe, und die Lehre herausnehme, so am verständlichsten und zur Besserung am meisten dienstlich. Wie ich denn hoffe, es sei hiemit treulich geschehen.

Daß ich aber des Ehrwürdigen Herrn, und meines lieben Präceptors, Doct. Luthers Summaria über den Psalter, habe darunter gesetzt, ist der Ursache geschehen, daß ich wohl wußte, ich würde es nicht besser machen, und sind dieselben bereits unter den Leuten gewesen. Wo dergleichen Summaria über die ganze Bibel vorhanden gewesen, wollte ich dieser Arbeit gar \unsicher{fein} müßig gestanden, und mich an dem, daß andere gemacht haben, gern begnügen lassen. Weil aber dergleichen nicht vor der Hand war, achte ich, es werde niemand solches mein Fürnehmen übel deuten noch strafen können, sintemal ich anders nichts, denn Erweiterung des Worts und Besserung der Kirchen damit gesucht habe. So ist es auch niemand gewehret noch verboten, der es kann, daß ers besser mache. Ohne allein, daß ich wohl weiß und sehe, daß Gott die muthwillige Verachtung und schreckliche Lästerung seines Worts an den Papisten dermaßen strafet, daß sie schier, wie die Juden, mit sehenden Augen starr blind sind, und Gottes Wort ihnen zum lautern Strick gerathen ist, wie man an ihrem Gemäsch und Lästern täglich sehet. Derhalb nicht zu besorgen, daß sie in diesem oder andern etwas verbessern. Besorgen muß man, wie sie von Tag zu Tag je länger je ärger und verstockter werden, also werde der Satan sie treiben, je länger je mehr Muthwillen zu üben, mit Verderbung ihrer selbst, und anderer, die Gottes Wort geringer achten, denn des Pabstes Geifer und Menschensatzung.

Der ewige barmherzige Gott, der solches Licht seines Worts uns so reichlich läßt scheinen, wolle mit Gnaden uns dabei erhalten, und dem Satan, der beide, durch den Türken und den Pabst, seine willigen Werkzeuge, damit umgehet, wie ers dämpfen oder fälschen möge, mit Ernst wehren, um seines lieben Sohnes und unsers Herrn Christi Jesu willen, durch seinen heiligen Geist; dem sei Lob und Ehre in Ewigkeit, Amen.

\quellseite{25}

\section{Widmung an Albrecht, Markgraf zu Brandenburg}

Dem Durchläuchtigen und hochgebornen Fürsten und Herrn, Herrn Albrechten, Markgrafen zu Brandenburg, in Preußen, zu Stettin, Pommern, der Caßuben und Wenden Herzogen, Burggrafen zu Nürnberg und Fürsten zu Rügen etc. etc.

Meinem gnädigsten Herrn:

Gnade, Friede und Barmherzigkeit, von Gott unserm Vater und unserm Herrn Jesu Christo.

Es ist eine sehr nütze und feine Ordnung jetzund hin und wieder in den Kirchen, so das Evangelium Christi angenommen, und von dem Gräuel des Papstthums abgetreten sind, daß man neben dem Predigen, alle Tage etliche Kapitel aus dem Alten und Neuen Testament von Wort zu Wort dem Volk vorlieset. Denn es will den Christen hoch vonnöthen sein, daß sie in steter Uebung des Worts bleiben, und nicht allein in der Historien der heiligen Schrift läufig und bekannt, sondern auch in den Predigten der Propheten erfahren sein.

Also siehet man, daß es im Judenthum eine sonderliche Ordnung gehabt hat, daß die Bücher Mose alle Sabbather für der Gemeine sind gelesen worden. Item, das Exempel Christi im Tempel, Luc. 4., zeigt auch an, daß man in den Synagogen der Propheten Bücher gelesen, und die Juden davon disputiret haben. Und der Apostel Sanct Paulus vermahnet seine Colosser, Cap. 3., daß sie das Wort Christi reichlich unter ihnen sollen wohnen lassen, in aller Weisheit, daß ein jeglicher sich selbst dadurch lehren und vermahnen möge.

Warum aber solches alles unter den Juden erstlich also geordnet, und hernach in der christlichen Kirchen auch gehalten, und heutiges Tages wieder in Brauch geführet ist, kann jedermann leichtlich urtheilen, der nur ein wenig geistlicher Erfahrung hat. Denn weil der Anfechtungen sehr viel sind, so ohne Unterlaß vom Teufel, der Welt, und unserm eigenen Fleisch erreget werden, ists unmöglich, daß man sich dagegen aufhalten, wehren und trösten möge, denn allein durch das Wort Gottes.

Darum siehet man auch, daß in der Schrift allenthalben Gottes Wort so herrlich gepreiset wird. Und sonderlich der 119. Psalm, neben andern schönen Namen, die er dem Wort Gottes gibt, nennet ers auch ein Licht, nach dem man gehen, und im Finstern sich für dem Anstoßen bewahren soll. Und hat hernach der heilige Apostel S. Petrus, 2 Petri 1., solchen Namen auch gebraucht. Denn weil wir hie auf Erden, unsers verderbten und sündlichen Fleisches, und anderer Aergerniß halben (gleichwie im Finstern) wandeln, und es sehr leicht geschehen kann, daß wir in Sünde oder Irrthum, entweder selbst fallen oder geführet werden, ist kein anderer Befehl, Rath noch Mittel, wider solche Fahr, denn daß die Christen dies Licht stets in Händen tragen, und nach demselbigen nicht allein ihre Wege und Wandel, so viel das ganze Leben betrifft, richten, sondern auch den Glauben rein und lauter, wider alle irrige Meinung erhalten.

Nun ist aber im Alten Testament viel Ding, in welche der gemeine Mann sich nicht richten kann, es sei denn, daß er zuvor eine Anleitung habe. Sonderlich aber siehet mans in den Propheten, daß unterweilen auch gelehrte Leute zu schaffen haben, daß sie die Predigt recht in einander richten, und die rechte Meinung treffen mögen. Denn der heilige Geist, ob er wohl durch die Propheten alles was von Christo und der Seligkeit zu wissen vonnöthen ist, genugsam hat angezeigt, so hat ers doch verborgener heimlicher Weise gethan.

Denn dies ist fürnehmlich unsers Herrn Christi Amt, daß er durch sein Evangelium vom Willen Gottes, Vergebung der Sünden, und unsrer Seelen Seligkeit uns lauter und klar, ohne alles Verdecken, unterrichten soll. Wie der heilige Johannes der Täufer von ihm sagt: Niemand hat Gott je gesehen. Der eingeborne Sohn, der in des Vaters Schooß ist, der hats uns verkündiget.

\quellseite{26}

Weil nun solches in den Propheten und dem Alten Testament dunkel ist angezeigt, und auch sonst viel Historien vorfallen, aus welchen sich nicht jedermann genugsam kann richten, ja da auch Aergerniß bei zu besorgen, wo man durch rechten Unterricht nicht zuvorkommt, habe ich meiner Kirchen hie zu gut, solche Arbeit auf mich genommen, und in das Alte Testament, auf alle Kapitel kurze Summaria gestellet, auf daß der gemeine Mann solche gewöhnliche Lection mit mehr Frucht hören, und allzeit eine kurze Lehre aus den Historien oder Predigten der Propheten fassen könnte, welche zur Besserung des Lebens und Unterricht des Glaubens nütz und dienstlich wäre.

Nun wirds nicht fehlen, solche Arbeit, so ich in meiner Kirchen hie zum besten habe vorgenommen, werde etlichen aus der Ursache nicht am besten gefallen, daß solche Summaria nicht alles fassen, das in den Historien und in der Propheten Predigten, in einem jeglichen Kapitel begriffen ist. Weil ich aber vornehmlich dahin habe müssen sehen, wie es solchem Kirchengebrauch am gelegensten wollte sein, hat es sich in keinen Weg wollen schicken, weil die Kapitel sonst lang genug sind, daß man allerwege eine lange Predigt und vollkommene Auslegung wollte daran hängen.

Habe mir also an dem genügen lassen, daß allerwege aufs kürzeste das angezeigt würde, welches zu merken am nöthigsten und nützlichsten ist. Denn auch neben dem, daß man solche Lectionen nicht soll zu lang machen, ists auch für den gemeinen Mann nicht, daß man alles weitläufig handeln, und die Kapitel in viel Stücke theilen wollte. Es ist genug, wo mit kurzen Worten allerwege das Fürnehmste angezeigt wird. Denn die vollkommene Erklärung gehört in die Predigt, nicht in solche Lection.

Es könnte auch wohl geschehen, daß man von etlichen Sprüchen und Historien schärfer und subtiler disputiren möchte, denn ich in diesen Summarien habe wollen thun, ich habe aber mich durch und durch beflißen, daß ich allenthalben dem jungen und einfältigen Volk zu gut, die gemeine und richtige Meinung anzeigte, die ihnen sein verständlich und nützlich wäre. Denn subtil disputiren, zu dem, daß es für den gemeinen Mann nicht taugt, kanns auch wohl in etliche Wege schädlich sein.

Um solche und andere dergleichen Ursachen mehr habe ich oft die Kapitel von einander getheilet, und zwei Summaria darauf gestellet, nicht allein darum, daß die Kapitel an ihnen selbst zu lang sind, wo man eine Ordnung mit dem Lesen hält, wie hie zu Nürnberg, sondern auch, daß oft die Händel darnach sind, daß es wohl bedarf, daß der Zuhörer ein jeglich Stück insonderheit bedenke, was es auf ihm habe, und wie ers ihm zur Besserung brauchen soll.

Item, ich habe auch viel Kapitel ausgelassen, und keine Summaria darüber gestellet, um der Ursachen willen, daß sie meines Erachtens, ohne sonderliche Frucht in der Kirchen heutiges Tages gelesen werden. Als da Mose den Tabernakel oder Hütten, mit seinem Zubehör, Salomon den Tempel und des Königs Haus bauet. Item, da Mose lehret, wie man mit dem Opfern umgehen, und die hohen Feste halten soll. Item, da man die Stämme des Volks Israel zählet rc. Denn ob man gleich solches in eine Allegorien ziehen, und eine heimliche Deutung daraus flechten will, wie es sich oft an etlichen Orten fein reimen läßt, so ists doch ohne Frucht, daß man öffentlich in der Kirchen den bloßen Text lieset.

Sonderlich aber halte ich keineswegs mit den Wiederläufern, die ihnen träumen lassen, die höchste Weisheit sei im Bau des Tabernakels und Tempels heimlich angezeigt. Eben wie die Juden auch narren, und halten das erste und zehnte Kapitel des Propheten Hesekiels für den heimlichsten Schatz der ganzen heiligen Schrift.

Wir Christen sollen das für das größte, und (wie es S. Paulus nennet) für das Geheimniß achten, so von Anbeginn der Welt ist verborgen gewesen, daß wir durch den Sohn Gottes vom Tod und der Sünden am Kreuz erlöset, und aus Gnaden, um Christi Willen, selig sind worden. Nach solchem Geheimniß ist das größte, daß Gott seinen Willen in den zehn Geboten offenbaret hat. Bei dieser Weisheit und Geheimniß sollen wir bleiben, und solches für das Höchste achten, wiederum alles anders, was in der heiligen Schrift und sonst ist, dem nirgend gleich rechnen. Weil nun das erste und zehnte Kapitel Hesekiels, ob es gleich eine nützliche Lehre in sich hat, dennoch so dunkel ist, daß man sich nicht wohl durchaus recht darein kann schicken, habe ichs in seinem Werth bleiben, und unter andern lassen liegen, welche ich für unnütz-

\quellseite{27}

thig achte, daß mans in der Kirchen lesen soll. Wer aber die Deutung solches Gesichtes begehrt zu wissen, kann es in der Vorrede Dr. Luthers über den Ezechiel lesen.

Nun bin ich Anfangs nicht Willens gewesen, solche Summaria in Druck zu geben, habe es auch eine zeitlang in meiner Kirchen allein brauchen lassen, aber weil sie doch endlich in fremde Hände kommen, und ich von vielen darum bin angesucht worden, daß ichs in Druck wollte geben, sintemal sie nicht allein dem gemeinen Mann, sondern auch vielleicht den schlechten einfältigen Kirchendienern, hin und wieder auf dem Lande, möchten zu Nutz kommen, habe ichs in die Länge nicht können weigern, ob, wie Hoffnung ist, der christlichen Kirchen ein wenig damit möchte gedienet werden.

Ich habe aber, Gnädigster Fürst und Herr, solche Arbeit fürnehmlich E. F. G. darum wollen dediciren, daß ich weiß, daß E. F. G. ein sonderlich Herz und Liebe zum Wort hat, und viel darauf gewagt, daß die Kirchen in E. F. G. Landen göttlich und recht würden angericht't. Denn gewiß ist es, wo E. F. G. nicht so treulich dazu geholfen und gerathen hätte, es wäre ganz Preußen noch heutiges Tages unter dem Pabstthum, und müßte des Evangelii eben sowohl, als die anstoßenden Herrschaften beraubt sein.

Daß aber nun GOttes Wort allenthalben in E. F. G. Landen recht geprediget, die Sacramente nach der Einsetzung Christi, den Leuten zum Trost und Seligkeit, recht gehandelt werden und GOttlob die Frucht des Worts nicht außen bleibt, sondern alle Tage viel Christen werden, und in Christo selig entschlafen, solches fürstlichen, ja bischöflichen und besten Werks ist E. F. G. durch Gottes Hülfe eine sonderliche Ursache.

Darum auch alle Christen E. F. G. billig lieb und werth haben, und für einen christlichen frommen Fürsten mit aller Wahrheit rühmen und loben, dagegen die andern, so um ihres Bauchs willen dem Pabstthum anhangen, und das heilige Evangelium, wo sie können, unbillig verfolgen, wohl werth sind, daß sie alle Welt hasse, und von ihnen sage, wie sie es mit ihren gottlosen Werken beweisen.

Sonderlich aber, gnädigster Fürst und Herr, hat mich zu solchem auch dieß verursacht, daß E. F. G. neulich mich ganz gnädigst hat lassen ersuchen, ob ich mich bei E. F. G. zu Kirchendiensten begeben wollte. Weil ich aber vorhin mit Diensten hie in meinem Vaterlande verpflichtet, und das Predigtamt sonderlich dieß auf ihm hat, daß niemand seinen Beruf ehe soll noch kann wechseln, denn man wolle sein nicht, habe ich mit solchem Dediciren E. F. G. meinen schuldigen Dienst und geneigten Willen wollen anzeigen, ob damit E. F. G. Kirchen, weil ichs gegenwärtig mit dem Leibe nicht kann, etwas möchte gedienet werden. Bitte derhalben E. F. G. wolle solches im besten annehmen, und mich E. F. G. gnädig lassen befohlen sein. Gott bewahre E. F. G. zu seinen Ehren, zu E. F. G. Seligkeit, und des ganzen Landes Wohlfahrt, Amen. Datum Nürnberg, am Tage Simonis und Judä, 1540.

Euer Fürstl. Gnaden

\begin{unterschrift}

unterthäniger und williger\\Vitus Dieterich,\\zu Nürnberg Prediger in der \unsicher{Sebalder} Pfarrkirchen.

\end{unterschrift}

\quellseite{28}

\begin{titelseite}

Vorreden und Beschlüsse\\[0.35em]

über die Kapitel der Bücher des Alten und Neuen Testaments,\\[0.35em]

wie dieselbigen sammt den Summarien Herrn Veit Dieterichs seligen, in der Kirche Gottes zu Breslau in den gewöhnlichen Kirchenlectionen, die Werktage über, zum Frühgebet, verlesen werden.\\[0.35em]

Gesetzt, und jezund neulich in Druck gegeben durch,\\[0.35em]

Franciscum Bierling,\\[0.35em]

Diener des göttlichen Worts daselbst, bei S. Maria Magdalena.\\[0.35em]

CUM PRIVILEGIO SENATUS.\\[0.35em]

ANNO 1676.

\end{titelseite}

\chapter{Widmung an den Rath der Stadt Breslau}

Den Edelen, Gestrengen, Ehrenvesten, Hochweisen und Wohlbenamten Herrn

J. Rathmannen

der Kaiserl. und Königl. Stadt Breslau, haltend und verwaltend die Hauptmannschaft Breslauischen Fürstenthums, Neumärkischen und Namslischen Weichbilder,

Meinen Großgünstigen Herren, Patronen, Förderern:

Gnade, Friede und Segen, von Gott dem Vater, durch seinen liebsten Sohn, unsern Herrn Jesum Christum, neben kräftigem Trost des heil. Geistes, sammt meinem andächtigen Gebet und unterthänigen Diensten, mit Fleiß zuvor.

Edle, gestrenge, ehrenveste, hochweise und wohlbenamte großgünstige Herren, es sind nunmehr acht und zwanzig Jahr verlaufen (dafür dem Vater aller Gnade und Barmherzigkeit, GOtt unserm Herrn, von dem alle gute Gabe und alle vollkommene Gabe von oben herab kommt, in wahrer Demuth, von ganzem Herzen, Lob, Ehre, Preis und Dank gesagt sei, durch JEsum Christum, seinen Sohn und unsern Mittler und Fürsprecher), daß ein ehrbarer, hochweiser und wohlbenamter Rath, bei den dreien Pfarrkirchen dieser löblichen Stadt anordnen hat lassen, daß neben den gemeinen öffentlichen Predigten auf die andern ledigen Tage der Wochen soll gehalten werden die Lection oder das Vorlesen der heiligen Biblia, der Bücher des Alten und Neuen Testaments, von den hochgelehrten Propheten und auserwählten Aposteln geschrieben, die von der christlichen Kirche Canonici libri, das ist, bewährte Bücher genennet werden, darum, daß die allein zu Zeugen, die Lehre der christlichen Religion zu bestätigen, sind zugelassen worden. Und dieses mit den Summarien des ehrwürdigen und hochgelehrten Herrn Magisters Veit Dieterichs (in GOtt seligen), welcher eine zeitlang ein fleißiger Schüler des werthen Mannes GOttes, Herrn Dr. Martin Luthers, und hernach Prediger zu Nürnberg in St. Sebalds-Kirche gewesen. Darinnen aufs kürzeste angezeigt wird, was am nöthigsten und nützlichsten ist dem jungen Volk und gemeinen Manne, aus allen Capiteln gemeldeter biblischer Schriften, zu wissen und zu lernen, darnach sie ihr Leben richten und solche seine Lehre zu ihrer Seelen Seligkeit brauchen können.

Diese nöthige und gute Ordnung hat ein ehrbarer, hochweiser und wohlbenamter Rath, als eine christliche Obrigkeit, die ihrer treuen Unterthanen zeitliche und ewige Wohlfahrt gern

\quellseite{29}

befördern wollte, billig machen lassen, und ist freilich auch alles Lobes werth, und als eine besondere große Gnade und Wohlthat GOttes zu erkennen, und mit Dankbarkeit allzeit zu preisen.

Denn es soll ja die theure, werthe, liebe heilige Biblia, und also die göttlichen Schriften der heiligen Propheten und Aposteln, dem Volke Gottes, sonderlich in den gewöhnlichen ordentlichen Kirchenversammlungen, täglich und fleißig vorgelesen werden.

Erstlich darum, daß GOtt, unser Herr, solches ernstlich befohlen hat, als der am besten verstehet, was uns Menschen und seinen Kindern auf Erden gut, nöthig und nütze ist. Gottes, des Vaters, Befehl ist, der also im 31. Kapitel des fünften Buches Mose aufgeschrieben stehet: Und Mose schrieb dies Gesetz und gabs den Priestern, den Kindern Levi, die die Lade des Bundes des HERRN trugen, und allen Ältesten Israel. Und gebot ihnen, und sprach: Je über sieben Jahr, zur Zeit des Erlaßjahrs, am Fest der Laubhütten, wenn das ganze Israel kommet zu erscheinen vor dem HErrn, deinem Gott, an dem Ort, den er erwählen wird, sollst du dies Gesetz vor dem ganzen Israel ausrufen lassen vor ihren Ohren, nämlich vor der Versammlung des Volks, beide der Männer und Weiber, Kinder und deines Fremdlings, der in deinem Thor ist; auf daß sie hören und lernen, damit sie den HErrn, ihren Gott, fürchten, und halten, daß sie thun alle Worte dieses Gesetzes, und daß ihre Kinder, die es nicht wissen, auch hören und lernen, damit sie den HErrn, euren Gott, fürchten, alle eure Lebtage. Diese Worte lehren, daß es Gottes endlicher und ernstlicher Wille sei, daß man die göttliche Lehre in den Büchern, welche man die heilige Schrift und Biblia nennet, verfasset, dem christlichen Volk fleißig vorlesen, bei demselben treiben, lehren, wiederholen und erklären soll, welches nicht allein in sieben Jahren einmal, sondern für und für alle Jahre und Tage im Schwange gehen soll; insonderheit aber auf alle Sabbathe, Sonntage und Feiertage. Der Sohn Gottes, unser HErr JEsus Christus, hat auch solches befohlen und geboten; denn also sagt er, Luc. 16.: Sie haben Mosen und die Propheten, laß sie dieselbigen hören. Und abermals Johannes im 5. Kapitel: Suchet in der Schrift (vernehmt durch stetiges Lesen, andächtiges Anhören und fleißige Betrachtung), denn diese ists, die von mir zeuget. Diese Mandate und Befehle unseres HErrn Christi soll ihm ein jeder Christ also einbilden, als rufe er ihm diese alle Tage zu, und soll derwegen keinen Tag dahin gehen lassen (sofern er nicht durch wichtige und nöthige Geschäfte verhindert wird), daß er nicht Gottes Wort sollte anhören. GOtt, der heilige Geist, befiehlet und vermahnet das auch durch Paulus, den Apostel und theuren Rüstzeug Gottes, Col. 3. und sagt: Lasset das Wort Christi unter euch reichlich wohnen, in aller Weisheit. Diese Worte sind ein solch göttliches Gebot, daß Gottes Wort bei uns soll in vollem Schwange gehen, und uns gar wohl bekannt sein, auch immer für und für bleiben. Denn wohnen ist an einem Ort eine gute lange Zeit sein und bleiben, da man nicht bald wieder abzieht. Es soll uns die heilige Biblia gar wohl bekannt sein, nicht anders, als diejenigen Leute, so bei uns und in einem Hause, oder ja hart neben uns wohnen, mit welchen wir täglich viel umgehen, Gespräche halten und zu thun haben. Es soll uns Gottes Wort also gemein sein, wie das tägliche Brod, welches wir alle Tage müssen haben und genießen, und dessen wir keinen Tag zur Nothdurft können entbehren. Denn wie das tägliche Brod eine leibliche Speise ist, die den Leib erhält und ernähret: also ist Gottes Wort eine geistliche Speise der Seelen, die uns muß erhalten zum ewigen Leben. Daraus folget, daß wir abermals einen starken Befehl haben, daß wir derwegen täglich Gottes Wort lesen oder anhören sollen.

Die andere Ursache, warum alle Tage Gottes Wort in den Kirchenversammlungen soll vorgelesen werden, ist die Nothdurft, daß es alle Menschen, jung und alt, in allen Ständen wohl bedürfen. Denn die Menschen sind von Natur zur Sünde nur sehr geneigt. So setzt der Teufel täglich an die Menschen mit sonderlichen Tücken und Stücken, daß er sie, zu ihrem Verderben, in Sünde und Schande locke, führe und stürze. Die Aergernisse der Welt-Kinder sind auch groß und überall gemein. Damit nun jedermann wider solche Gefahren und Unfälle verwahret werde, ists ja nöthig, daß täglich zur Warnung Gottes Wort vorgelesen werde, damit sich niemand sein sündlich Fleisch und Blut, den bösen neidischen Teufel und die ärgerliche Welt in Sünden an-

\quellseite{30}

führen und zu seinem zeitlichen und ewigen Schaden verführen lasse. Denn das göttliche Wort muß vorleuchten in allem unserm Vornehmen, Rath und That, Thun und Lassen, sollen wir anders unser ganzes innerliches und äußerliches Leben und Wandel recht anstellen; auch unsere bösen Affecten und Zuneigungen dämpfen und im Zaum halten, und unserer Vocation und Beruf treulich vorstehen wider alle Aergernisse und Hindernisse. Denn die heilige Schrift ist die einige Norm und Form, Regel und Richtschnur, nicht allein des rechten seligmachenden Glaubens, sondern auch aller Gottseligkeit, als die da lehret, wie man Gott in rechtem Glauben und gutem Gewissen, wahrhaftig und gefällig dienen soll. Solches bezeuget der königliche Prophet David in Frage und Antwort, da er spricht im 119. Psalm: Wie wird ein Jüngling des Herrn Weg unsträflich gehen? Wenn er sich hält nach seinen Worten und Geboten. Und sagt ferner: Ich behalte dein Wort in meinem Herzen, auf daß ich nicht wider dich sündige. Dein Wort ist meines Fußes Leuchte, und ein Licht auf meinem Wege. So ist auch Gottes Wort eine starke Verwahrung und Antidotum wider die Sünde, wider die Anreizung des Teufels und wider die Aergernisse der gottlosen Epicurer und Sünder, die sie den Christen an den Weg, darüber sie gehen müssen, legen und eingraben. So sind auch unter den wahren Christen allzeit betrübte Herzen, die das Joch des heiligen Kreuzes tragen, jetzt am Leibe, jetzt an ihrer Nahrung, in diesen ganz kümmerlichen und geschwinden Zeiten, bald an ihren Ehegatten und Kindern, und die mit Schwermuth und Anfechtung beladen sind, daß sie in ihrem Nothstande unter Gottes gewaltige Hand sich demüthigen, und solch ihr Kreuz sich nicht unterdrücken lassen, sondern täglichen nöthigen Trost haben, mit Geduld ihre Seele fassen und in wahrem Glauben der verheißenen gewissen Hülfe Gottes erwarten. Derwegen erfordert's abermals die Noth solcher und anderer Kreuzträger und Kreuzträgerinnen, daß ihnen zu Trost Gottes Wort täglich vorgelesen werde. Denn die heilige Schrift ist auch das rechte Trostbuch, daraus allein angefochtene, bekümmerte und geängstigte Herzen zu jeder Zeit rechten, gewissen und beständigen Trost schöpfen und erlangen mögen, durch Gnade des heiligen Geistes, deß Buch es ist, als des rechten himmlischen Kanzlers. Davon zeugt abermals David, der israelitische König und Prophet Gottes im 119. Psalm und sagt: Wenn dein Gesetz nicht wäre mein Trost gewesen, so wäre ich vergangen in meinem Elende. Und der heilige Prophet Jeremias im 15. Kapitel danket dem Herrn für solche Kraft und Nutzbarkeit seines göttlichen Wortes, und spricht ihm also zu: Indeß enthalte uns dein Wort, wenn wirs kriegen; und dasselbe dein Wort ist unsers Herzens Freude und Trost. Von des göttlichen Wortes kräftiger und heilsamer Wirkung zeuget auch der heilige Apostel Paulus im 15. Kapitel seiner Epistel an die Römer, mit davon folgendem Lobsspruch: Was zuvor geschrieben, das ist uns zur Lehre geschrieben, auf daß wir durch Geduld und Trost der Schrift Hoffnung haben. Ferner ist auch dies eine große Nothdurft aller Menschen, weil wir leben und wissen nicht, wie lange, und weil wir sterben, und wissen nicht, wann; daß ein jeder, wie mit dem Namen, also auch mit der That, ein wahrer Christ sei und bleibe, da er anders nicht soll in Sünden sterben und ewig verderben. Also ist auch die Frage, wer denn ein rechter Christ sei? Antwort: Ein rechter Christ ist der Mensch, welcher in täglicher wahrer christlicher Buße lebet; der alle Tage Gott, den himmlischen Vater, im Namen Jesu Christi, seines lieben Sohnes, bittet um Vergebung aller seiner Sünden; der auch ein gläubig Herz in seinem Leibe augenblicklich herumträgt, daß er mit Gott versöhnet sei, und habe Gnade erlanget um des Herrn Christi Gehorsams und Leidens willen; der fleißig betet für sich, die Seinen und für die gemeine Noth der Christenheit; der unter dem von Gott auferlegten Kreuz geduldig ist; der in seinem Beruf vor wissentlichen, muthwilligen und vorsätzlichen Sünden sich hütet, und dagegen züchtig gegen ihn selbst, gerecht gegen seinen Nächsten und gottseelig gegen Gott lebet, so viel ihm möglich, durch Hülfe des heiligen Geistes, und der, da er aus menschlicher Schwachheit oder Verführung der Welt und des Teufels einen Fehltritt thut, mit bösen Gedanken, Worten und Werken sündiget, bald solches in wahrem Glauben Gott abbittet, und der im Herrn Christo endlich selig entschläft und aus dieser Welt scheidet. Diese Bewegungen, Werke und Tugenden des wahren Christenthums kommen aus Gottes Wort, und werden dadurch gestärket und vermehret von Tag zu Tag in denen, die Gottes Wort mit Ernst, Andacht und Fleiß hören, und im Gehorsam desselbigen sich üben. So erfordert's demnach wiederum die höchste Nothdurft aller Menschen, daß sie aus Gottes Wort davon gelehret und dazu vermahnet werden.

\quellseite{31}

Daß aber die heilige Schrift und Gottes Wort also von wahrem Christenthum lehre und unterrichte, bekräftigen folgende Zeugnisse, insonderheit unsers, der Heiden, Apostels S. Pauli:

Röm. 3. sagt er: Durch das Gesetz kommt Erkenntniß der Sünden.

Röm. 10.: So kommt der Glaube aus der Predigt, das Predigen aber durch das Wort Gottes.

2 Cor. 5.: Gott hat uns mit ihm selber versöhnet durch Jesum Christ, und das Amt gegeben, das die Versöhnung prediget. Denn GOtt war in Christo, und versöhnte die Welt mit ihm selber, und rechnete ihnen ihre Sünde nicht zu, und hat unter uns aufgerichtet das Wort von der Versöhnung.

Röm. 10.: Wer den Namen des HErrn wird anrufen, soll selig werden. Wie sollen sie aber anrufen, an den sie nicht glauben? Wie sollen sie aber glauben, von dem sie nichts gehört haben? Wie sollen sie aber hören ohne Prediger?

Röm. 15.: (Damit dieser nöthige Spruch abermal wiederholet werde). Auf daß wir durch Geduld und Trost der Schrift Hoffnung haben.

2 Tim. 3.: Alle Schrift von Gott eingegeben, ist nütz zur Lehre, zur Strafe, zur Besserung, zur Züchtigung in der Gerechtigkeit, daß ein Mensch Gottes sei vollkommen, zu allem guten Werk geschickt.

Röm. 1.: Das Evangelium von Christo ist eine Kraft Gottes, die da selig macht alle die daran glauben. Und 1 Cor. 1.: Dieweil die Welt durch ihre Weisheit Gott nicht erkannte, gefiel es Gott wohl, durch thörichte Predigt selig zu machen die, so daran glauben.

Die dritte Ursach, daß man auch heute noch Gottes Wort täglich vorlesen, und sich darin üben solle, sollen sein die löblichen Ordnungen und besondern Exempel der Heiligen Gottes, und theuren werthen Männer in der heiligen christlichen Kirchen, als:

Im 12. Kapitel des 1. Buchs Mose wird gesagt von dem Patriarchen und Altvater Abraham, daß er eine Hütten aufgerichtet hat, und bauete dem HErrn einen Altar, und predigte vom Namen des HErrn: Das ist, er hat die Seinen alle Tage gelehret von dem verheißenen Messias, und sie erinnert seiner Wohlthaten, wie solches zu dieser Zeit geschieht durch das tägliche Vorlesen, und die Predigten göttlichen Worts.

Das 13. Kapitel der Geschichte der Apostel meldet deutlich, daß im Judenthum eine sonderliche Ordnung gewesen, daß das Gesetz oder die Bücher Mosis und der Propheten alle Sabbathter in der Synagogen und Judenschulen sind vorgelesen worden.

Dieses bestätiget auch das Exempel unsers HErrn Jesu Christi, Luc. 4., denn da er gen Nazareth kömmt, da er empfangen und erzogen war, und am Sabbath in die Schule ging, und daselbst nach seiner Gewohnheit aufstehet und lesen will, da wird ihm das Buch des Propheten Jesaia gereicht, daraus er hat eine tröstliche Predigt gethan, von der gegenwärtigen Gnadenzeit, davon der gemeldete Prophet im 61. Kapitel geweissaget hat.

In Geschichten der heil. Apostel im 15. Kapitel wird abermal von solcher jüdischen Ordnung Zeugniß gegeben. Denn also spricht S. Jacobus: Moses hat von langen Zeiten her in allen Städten, die ihn predigen, und wird alle Sabbathter in den Schulen gelesen. Und zuvor im 13. Kapitel wird auch gedacht der Lectio des Gesetzes und der Propheten, welche in den Judenschulen am Sabbath üblich, und in stetem Gebrauch gehalten worden sind. Aus diesem allen ists offenbar, daß, wie auch vor Christi Geburt, also dergleichen zu seiner Apostel Zeiten und Lehrtagen, bei dem Volk Gottes allwegen der Gebrauch gewesen, daß man das Gesetz Mosis, sowohl der heiligen Propheten Schriften, beides im Tempel zu Jerusalem und in allen Synagogen und Schulen, im ganzen Lande, nicht allein alle Feste und Sabbathte, sondern auch täglich dem Volk vorgelesen und ausgelegt hat, wie auch heutiges Tages bei den Judenschulen, wie sie jetzt geschaffen, zu sehen.

Diese nothwendige Uebung der heil. göttlichen Lehre haben auch der Apostel unsers HErrn Christi, treue Successoren und Nachkommen, die alten christlichen Kirchenlehrer, behalten, und haben die Schriften der heil. Propheten und Apostel in ihren Gotteshäusern und öffentlichen Ver-

\quellseite{32}

sammlungen, beides auf die Fest- und Feiertage, und sonst wo man zusammenkommen ist, der christlichen Gemeine und Versammlung vorgetragen, vorgelesen und ausgelegt. Solches bezeugen:

Justinus, der Lehrer und Märtyrer Christi, welcher gelebet hat um das 150. Jahr nach Christi Geburt, und sagt in der andern Verantwortung für die Christen: Am Sonntage kommen alle, die in Städten und auf dem Lande wohnen, zusammen an einem Ort, und werden der Propheten und Apostel Schriften gelesen, so viel daß eine Stunde leiden kann. Darnach, wenn der, so gelesen hat, aufhöret, erinnert und vermahnet der Vorsteher das Volk, daß sie das Gute, so vorgelesen, wollten nachthun.

Clemens Alexandrinus, der berühmt gewesen ist um das 200. Jahr, sagt im 7. Buch Stromatum, daß der Christen Opfer gewesen sein die Gebete und Lobgesänge, und die Lectionen der Schrift, so für Essens geschehen.

Mit diesen stimmet Tertullianus, der gelebt hat um das 230. Jahr, im 39. Kapitel seines Apologetici adversus Gentes. Wir kommen, spricht er, zusammen zur Predigt und Erklärung der Schrift, wenn uns die Gelegenheit der gegenwärtigen Zeit zwinget, etwas zu vermahnen oder zu erinnern. Wahrlich, wir speisen unsern Glauben mit dem heiligen Worte, wir richten die Hoffnung auf, wir stärken das Vertrauen, und mehren nichts desto weniger die Zucht mit stetigem Einbläuen der Gebote.

Mehr Zeugnisse der ersten Kirchen anzuziehen, ist nicht vonnöthen. Und weil Justinus in Asia, als nämlich in dem Städtlein Nicopoli, Palästina, geboren, desgleichen zu Rom und anderswo in Europa, gelebet hat; Tertullianus in Afrika; Clemens zu Alexandria in Egypten, ist daraus offenbar, daß zu der Zeit an allen Orten, so viel den Gebrauch und die Weise die heilige Bibel zu lesen anlanget, eine wunderliche und sonderliche Einigung gewesen, und Gleichförmigkeit gehalten worden sei.

Von diesem alten und löblichen Brauch haben unter andern Dienern der Kirchen die Lectores auch ihren Namen gehabt, von ihrem Amt, daß sie die heilige Schrift in der christlichen Versammlung haben gelesen und vorgetragen, welchen das Amt zu lesen vom Bischofe mit Uebergebung des heiligen Buchs der Biblia, in Gegenwärtigkeit der ganzen Gemeine befohlen und vertraut ward.

Hieher gehören die Exempel vortrefflicher und heiliger Leute, die fleißig in Gottes Wort sich geübet, und die heilige Biblia täglich gelesen haben.

David, der tapfere Held in Frieden und Streit, der herrliche König in Krieg und Sieg, da er jetzt vor dem Tyrannen Saul flohe, und in unwegsamen Orten in der Irre ging, jetzt beim Könige Gath im traurigen Elende war, oder als er sich für seinem Sohne Absalom, der einen bürgerlichen Aufruhr wider ihn erwecket, fürchtete, hernach, als er draußen im Felde stritt oder daheim das Regiment versahe, war er dennoch auf das Gesetz und Gottes Wort, welches er auf Gottes Befehl angeschrieben und allzeit bei sich hatte, beflissen, ja auch in den alltraurigsten Empörungen und Unruhen hat er selbst sehr viele Psalmen geschrieben und gesungen.

Dergleichen haben gethan der mühselige Salomon, der geschäftige Hiskia, der überladene Josias, die Könige in Juda.

Der edle Kämmerling oder Schatzmeister der Mohren-Königin Candaces, zu welchem, wie glaublich ist, auch das Evangelium kaum von Christo kommen war, brachte die Zeit zu mit der Lection des Propheten Jesaia, nicht allein daheim, unangesehen, daß er beschweret genug und unmüßig war, wegen seines verdrießlichen Amts, belangend die Rechenschaft des Einnahmens und Ausgebens, sondern auch da er über Feld reisete, und auf dem Wagen geführet ward, in der Geschichte der Apostel im 8. Kapitel.

Vom Kaiser Theodosio Secundo melden die Historien, daß er mit eigener Hand die Bücher des Neuen Testaments abgeschrieben, und täglich zur Morgenzeit daheim ein Stück in demselben gelesen hat: Und daß ich dies auch melde, so hat er auch Psalmen und Gebetlein mit seiner Gemahl, Eudoxia genannt, und mit seinen Schwestern zu singen und zu recitiren oder zu sprechen gepflegt.

Vom Alfred, dem Könige aus England, welcher die berühmte Hochschule zu Oxford oder Ochsenfurt gestiftet hat, zeuget die Chronica, daß er die ganze Zeit des Tages in drei

\quellseite{33}

Theile getheilet habe, und täglich acht Stunden mit Lesen, Beten und geistlichen Betrachtungen zugebracht, und denn so viel mit der Regierung des Reichs. Die andern Stunden mit dem, was auf Pflegung des Leibes gehöret.

Eben auf diese Weise hat Carl der Große, der erste römische deutsche Kaiser, sein Gemüth geweidet mit dem Studiren, zuvor aber mit der Lection der heiligen Schrift. Auch hat er unter Essens allezeit etwas aus der heiligen Bibel, dazu auch aus den Schriften des Kirchenlehrers Augustini ihm lassen lesen.

Alphonsus, der hochweise und gelehrte, auch der Lehrer und Gelehrte liebhabende König in Sicilien und zu Neapolis, hat die heilige Biblia mit den Explicationibus oder Erklärungen Nicolai de Lyra, dreizehn mal ausgelesen.

Also haben fleißig die heilige Biblia selber gelesen, und haben ihnen und den Ihren lesen lassen, Herr Christianus III., der christliche löbliche König in Dänemark. Also Frau Margaretha, die heilige, fromme Fürstin zu Anhalt, geborne Schlesische Herzogin von Münsterberg, und andere viel mehr fürstliche und andere christliche Personen.

Aus diesen und dergleichen Exempeln folget, weil alle wahre Christen die Schriften der heiligen Propheten und Apostel herzlich gern täglich mit Andacht gelesen, und lesen haben lassen, daß also öffentlich im gemeinen die heilige Biblia auch billig alle Tage soll vorgelesen werden, und daß jedermann die Lection gern und mit allem Fleiß auch anhören, und derer abwarten soll.

Und dies sind allein drei fürnehme Ursachen, welche alle Menschen, so sich der christlichen Religion rühmen, mit Ernst, gleich als mit einer Posaune und Feldgeschrei, zur Anhörung der Lection der heiligen Biblia, aufmuntern, reizen und treiben sollen, also, daß sich auch niemand daran hindern lassen soll: weil keines Standes, keiner Verrichtung irgend eines Amts, keiner Arbeit, keiner vorfallenden Sache, von den gemeinen und eigenen Händeln, Entschuldigung billig genugsam sein kann, daß sie irgend einen Christen los mache vom Anhören oder ja Lesen der heiligen biblischen Bücher.

Weil nun aber erzähleter und anderen Ursachen wegen, auch ein ehrbarer, hochweiser, wohlbenamter Rath, als eine christliche und gottselige Obrigkeit, angeordnet hat die Lection der heiligen Biblia, ward rathsam, nöthig und nützlich befunden, damit zu fleißiger andächtiger Anhörung, und nöthigem heilsamen Gebrauch göttlichen Worts die lieben Zuhörer angebracht und verursacht würden, daß allezeit ein jedes Kapitel der heiligen Biblia mit einem sonderlichen Prolog oder mit einer Vorrede angefangen, und daß ein jedes Summarium des wohl gemeldeten seligen Herrn M. Dieterichs mit einem Epilogo, Voto, Wunsch oder Gebetlein sollte beschlossen werden. Dieses ins Werk zu richten, hat der ehrwürdige und hochgelehrte Herr Lucas Pollio, das theure Werkzeug Gottes und wohlverdiente Pfarrherr bei unserer Kirchen zu S. Maria Magdalena, damals mir vertraut und auferleget, daß ich solche Prologos oder Vorreden, desgleichen die Epiloge und Beschlüsse ordentlich stellen und fassen sollte. Solches ich auch im Namen Gottes auf mich genommen, und habe diese Arbeit mit Rath und Anweisung wohl gemeldeten meines seligen Pastors und lieben Mitbruders in Christo, auch von vielen Jahren treuen Freundes, und nun wegen seiner geliebten Schwester, meiner ehelichen Hausfrauen, Herrn Schwagers, also zusammengebracht und geschrieben.

Gleichwie aber auch der theure Mann Gottes, Herr Philippus Melanchthon (dessen Auditor und Schüler ich gewesen bin etliche Jahr, und habe seine heilige Cygneas Cantiones angehört, welches ich unter die vornehmsten und größten Wohlthaten rechne, die mir mein lieber GOTT bisher mein Lebenlang aus Gnaden und Barmherzigkeit erzeiget hat), alle seine Scholas, Lehre und Vermahnungen also und dahin richtet, auf daß er die wahre Gottesfurcht in uns, seine Discipul und Zuhörer, pflanzte, und uns zu täglichem, herzlichem, gläubigem Gebete gewöhnte, auch nöthigen Trost weisete, weil er allezeit sagte: wir würden betrübte gefährliche Zeiten und Läufte erleben; (wie denn seine Weissagung nicht ohne sonder große Schmerzen, Kreuz und Noth vieler frommer Herzen bisher erfüllet worden und noch erfüllet wird), also habe ich die gemeldeten Prologe gestellt, daß der meiste Theil Vermahnungen sind zu täglicher Buße, das ist, zur Furcht Gottes, zur Bekehrung, zum Glauben, Gebet und Uebungen der Tugenden im Decalogo oder zehen Geboten Gottes befohlen; viel Prologe und Vorreden habe ich auch accordirt und vereinigt mit den Summarien, daß sie in specie, gewisse nöthige

\quellseite{34}

Lehrpuncte des Texts und der Kapitel mit erklärten und zu besonderen Haupt-Tugenden reizen; damit abermal Gottesfurcht, Gebet und Trost in den Herzen der Zuhörer, der sichern und furchtsamen, der betrübten und fröhlichen erwecket, gepflanzet und erbauet würde. Damit auch die Lehrpuncte, wie die in den Summarien des Herrn Veit Dieterichs seligen berühret werden, mit gläubigem Herzen, und mit freudigem angenehmen Willen, die ausständige Zeit des Tages, und das ganze Leben über, im Gehorsam Gottes zu leben, haften und Nutzen bringen möchten, derwegen, wie alle ordentlichen Predigten mit einem christlichen Wunsch beschlossen worden, daß das angehörte Wort Gottes bleibe und bestehe in einem guten und feinen Herzen, und hundertfältige Früchte in Geduld und allen Gott wohlgefälligen Tugenden bringe, also habe ich auf die gemeldete Summaria ein Votum und Beschluß gestellet, zur Application, wie es dieselben erfordern, und der gütige Leser solches wohl erkennen wird können. Und habe mich beflißen, so viel möglich, sonderlich um der Einfältigen willen, daß ich gewöhnliche, deutliche und verständliche Worte gebraucht.

Erwähnte Arbeit ist also, durch Gottes väterlichen reichen Segen gerathen und angewandt worden, daß diese Lectionen nicht allein fleißig und andächtig unsere christliche Kirchenversammlung angehört, sondern daß auch viele von unserer gottseligen Bürgerschaft, nachdem der Geist Gottes in ihnen gewirket, einer jetzt des Stückes, bald ein anderer eines andern, Abschrift begehrt haben; daraus denn auch weiter folget, daß von etlichen Jahren her ist angehalten worden, daß solches alles in Druck verfertiget würde, damit solche nöthige und heilsame Lehren und Erinnerungen auch zu Hause den Kindern und dem Gesinde zu Nothdurft und Nutzen möchten vorgelesen werden.

Dieses (das den begehrten Abdruck anlanget), edle, gestrenge, ehrenveste, wohlweise, wohlbenamte großgünstige Herren, an eure Gnaden und Herrschaften billig, aus schuldiger Pflicht und Unterthänigkeit, gelangen hat sollen, und da solches in Gehorsam von mir geschehen, haben eure Gnaden und Herrschaften, auf darüber vor gehaltenen Rath, günstig bald bewilliget, daß solches Buch in dieser Stadt-Druckerei solle publiciret und verfertigt werden.

Daß nun solches also gerathen, danke ich von Herzen dem treuen lieben Gott, der es also gesegnet, daß es eure Gnaden und Herrschaften und vielen frommen Herzen gefallen; und lasse auch nun im Namen meines Herrn Christi diese Arbeit, so gut, als er sie gegeben, in offenem Druck ausgehen, auf daß also diese löbliche und vortreffliche Stadt und christliche Gemeine einen Dank von mir habe für die vielfältigen und großen Wohlthaten, die ich allhier empfangen habe: anfänglich in meiner Jugend, in der ich von der Neiße (die meine liebe Geburts-Stadt ist), hieher ad studia pietatis et bonarum artium kommen bin. Da ich denn bei ehrbaren und vornehmen Leuten ehrliche und gar bequeme Paedagogias gehabt, und denn darauf von einem ehrbaren, hochweisen wohlbenamten Rath, unangesehen, daß ich ein Fremder war, zu ihrer Gnaden und Herrschaften Alumno angenommen, und zu Wittenberg Verlag gehabt durch Beförderung des gestrengen, edlen, ehrenvesten und wohlbenamten Herrn Anton Band, Hauptmann etc., in Gott seligen, sowohl auch durch Hülfe des edlen, ehrenvesten und wohlbenamten Herrn Hans Morenberg, des Rathes, und Scholarchae oder Schul-Herrns, auch seligen, welches ich mit dankbarem Herzen und Gemüthe wohlgemeldeten Herren und ihren Erben zu Ehren billig rühme, mit herzlichem Wunsch, daß der barmherzige und gütige Gott diese Wohlthaten an wohlgemeldeten Herren, meinen großen Patronen und Förderern, ehrbaren Geschlechtern und Nachkommen, allen und jeden, allzeit, bis an den jüngsten Tag, reichlich, an Leib und Seele, zu zeitlicher und ewiger Wohlfahrt, vergelten wolle. Darnach, auf daß ich auch für die viele Ehre, Freundschaft und große Wohlthaten, so diese 31 Jahre, weil ich in dem ehrwürdigen Kirchenamt diene, mir, meinen lieben, den ersten und andern Eheweibern und Kindern, von vielen frommen Herzen erzeiget sind worden, mich etlicher Maßen dankbar erfinden lasse. Denn ich ja auch die wiederum ehren soll, welche mir Ehre und Gutes bewiesen haben. Wie solches die Vernunft lehret, daß es billig sei, und die heilige Schrift, daß es geschehen solle, ernstlich befiehlet, mit Ankündigung greulicher Strafen gegen die, so undankbar erfunden werden.

Diese Vorreden und Beschlüsse aber über den ersten Theil der heiligen Biblia, nämlich über das Alte Testament, haben eure Gestrengen und Herrschaften, mit schuldiger Reverenz und Ehrerbietung ich sollen und wollen dediciren und zuschreiben, damit ich für meine Person, Gott zu

\quellseite{35}

Ehren und Dank, öffentlich wahrhaftig Zeugniß mit gebe, daß eure Gnaden und Herrschaften auch rechte Pfleger und Vormünder, und treue Säugammen und Wärterinnen seien der Kirchen des Herrn Christi, so unter Euer Gnaden und Herrschaften löblichen Gemeine und lieben Unterthanen, ihm der Sohn Gottes sammlet; und daß Eure Gnaden und Herrschaften als christliche Herren und Regenten auch wahre Theophili und Gottes Liebhaber seien. Wie denn der heilige Evangelist Lucas sein Evangelium und Scripta auch seinem Theophilo dediciret und zugeschrieben hat.

Gelanget demnach an Eure Gnaden und Herrschaften sämmtlich und sonderlich, mit unterthänigem und demüthigem Fleiße, meine Bitte, Eure Gnaden und Herrschaften wollen ihnen diese meine wohlgemeinte Dedication, auch Declaration und Erzeigung meiner schuldigen Dankbarkeit günstig gefallen lassen, die mit gutem Willen annehmen; auch mein, meines lieben Eheweibes, ihres Sohnes und meiner Kinder, und sonderlich nach meinem tödtlichen Abgang, Vater und Förderer sein und bleiben, und also uns in ihre Gunst und Patronium, zu väterlicher Beförderung, lassen befohlen sein. Das wollen wir mit unserm andächtigen Pater noster und Gebet, für Euer Gnaden und Herrschaften und gemeiner Stadt Wohlfahrt, Gedeihen, zeitlichen und ewigen Segen treulich zu flehen, nach höchstem und besten Vermögen zu jeder Zeit beflissen sein.

Helffe der treue, barmherzige und gütige Gott, daß alle Menschen und Herzen, welche die Lection der heiligen Biblia, der Summarien darüber, und dieses Buches anhören werden, daß sie je länger je mehr erleuchtet werden, und daß sie ihres christlichen Berufes gewisser, in gottseligem Wandel fleißiger, zum Creuz williger, in Anfechtung getrösteter, zum Gebet ernster und brünstiger und zu allen guten Werken geschickter werden, und also zu Gottes Ehre, und allen, die solche Lection anhören werden, zur Seligkeit gedienet werde in Christo Jesu, unserm Herrn, bis zu dem gewünschten jüngsten Tage, Amen.

Auch bitte ich herzlich, der allmächtige Gott und treue Vater wolle ihm E. G. und Herrschaften, diese löbliche Stadt, Kirchen und Schulen und die ganze ehrbare Gemeine, in seine Gnade und Schutz lassen befohlen sein; wolle Eure Gnaden und Herrschaften lange Zeit frisch und gesund, bei glückseligem Regiment, gesegneter Haushaltung und aller Wohlfahrt Leibes und der Seelen, um seines geliebten Sohnes, unsers Herrn und Heilands Jesu Christi willen, für und für erhalten, schützen und bewahren, und nach vollendter Ritterschaft dieses elenden Lebens, die himmlische Ehren-Krone ewiger Freude und Seligkeit verleihen, Amen. Datum im Jahr Jesu Christi, unsers Herrn und Seligmachers Geburt 1596, den 18. Februar, an welchem Tage das auserwählte Rüstzeug Gottes, der Herr Dr. Martin Luther, seliger Gedächtniß, als der treue Gott sein vorgehabtes Werk mit großer Kraft und Stärke durch ihn ausgerichtet, zu Eisleben, in seiner Geburtsstadt, vor 50 Jahren in Christo, dem Herrn, fröhlich und selig entschlafen, seines Alters im 63. Jahr.

\begin{unterschrift}

E. G. und Herrschaften\\Unterthäniger Diaconus der Kirchen zu S. Maria Magdalena Franciscus Bierling.

\end{unterschrift}
